\documentclass[preprint,11pt,authoryear]{elsarticle}

%─── Packages ────────────────────────────────────────────────────────────────
\usepackage[T1]{fontenc}
\usepackage{lmodern}
\usepackage{amsmath,amssymb,amsthm,bm,bbm}
\usepackage[dvipsnames]{xcolor}
\usepackage{booktabs,colortbl,array,multirow,threeparttable}
\DeclareMathOperator*{\argmin}{arg\,min}
\usepackage{graphicx,float,caption,subcaption}
\graphicspath{{figures/}}
\usepackage{algorithm,algpseudocode}
\usepackage{natbib}
\usepackage{hyperref}
\hypersetup{colorlinks=true, allcolors=black, citecolor=black,
            linkcolor=black, urlcolor=black, pdfborder={0 0 0}}
\usepackage{placeins}
\usepackage{setspace,microtype,enumitem}
\onehalfspacing

%─── Generated figures ───────────────────────────────────────────────────────
% Every number the text asserts comes from scripts/paper_numbers.py, which
% recomputes it from the artefacts and fails the build if a value has drifted.
% Generated by scripts/paper_numbers.py -- do not edit by hand.
% Every figure the manuscript asserts, recomputed from the artefacts.
% Macro names carry their panel: Main = 16 x 24 = 384 pairs (Table 1);
% Seq = 13 x 24 = 312 pairs, the forecasters with stored series.

\newcommand{\nMainForecasters}{16}
\newcommand{\nMainAssets}{24}
\newcommand{\nMainPairs}{384}
\newcommand{\nSeqForecasters}{13}
\newcommand{\nSeqPairs}{312}
\newcommand{\nMainGreen}{335}
\newcommand{\nMainGreenPct}{87.2}
\newcommand{\nMainKupiecRawPasses}{123}
\newcommand{\nMainKupiecCorPasses}{273}
\newcommand{\nMainBestKupiec}{15}
\newcommand{\nMainBestKupiecModel}{CAViaR-AS}
\newcommand{\nMainCCDefined}{180}
\newcommand{\nMainCCUndefined}{204}
\newcommand{\nMainCCUndefPct}{53.1}
\newcommand{\nMainCCPass}{97}
\newcommand{\nMainCCPassPct}{53.9}
\newcommand{\nMainCCAsPassPct}{78.4}
\newcommand{\nMainRawPiMin}{0.0110}
\newcommand{\nMainRawPiMax}{0.0333}
\newcommand{\nMainRMin}{0.001}
\newcommand{\nMainRMax}{0.357}
\newcommand{\nMainRTruncOne}{17.3}
\newcommand{\nMainRTruncTwo}{23.5}
\newcommand{\nRawPiTimesFM}{0.0143}
\newcommand{\nRawPiMoiraiOne}{0.0154}
\newcommand{\nRawPiMoiraiTwo}{0.0178}
\newcommand{\nRawPiLagLlama}{0.0294}
\newcommand{\nRawPiChronosSmallA}{0.0173}
\newcommand{\nRawPiChronosMiniA}{0.0177}
\newcommand{\nRawPiGJR}{0.0200}
\newcommand{\nRawPiGJRt}{0.0145}
\newcommand{\nRawPiCAViaRAS}{0.0110}
\newcommand{\nWidthRatioTimesFM}{1.016}
\newcommand{\nWidthRatioMoiraiTwo}{1.001}
\newcommand{\nGJRnegqV}{0}
\newcommand{\nSeqKupiecRejRawOne}{69.9}
\newcommand{\nSeqCCUndefRawOne}{34.6}
\newcommand{\nSeqCCRejRawOne}{27.5}
\newcommand{\nSeqCCUndefCorOne}{53.5}
\newcommand{\nSeqCCRejCorOne}{48.3}
\newcommand{\nSeqKupiecRejRawTwoFive}{62.8}
\newcommand{\nSeqCCUndefRawTwoFive}{11.9}
\newcommand{\nSeqCCRejRawTwoFive}{19.3}
\newcommand{\nSeqCCUndefCorTwoFive}{26.6}
\newcommand{\nSeqCCRejCorTwoFive}{30.6}
\newcommand{\nSeqKupiecRejRawFive}{47.8}
\newcommand{\nSeqCCUndefRawFive}{1.0}
\newcommand{\nSeqCCRejRawFive}{22.0}
\newcommand{\nSeqCCUndefCorFive}{4.5}
\newcommand{\nSeqCCRejCorFive}{24.5}
\newcommand{\nSeqKupiecRejRawTen}{40.1}
\newcommand{\nSeqCCUndefRawTen}{0.0}
\newcommand{\nSeqCCRejRawTen}{26.0}
\newcommand{\nSeqCCUndefCorTen}{0.0}
\newcommand{\nSeqCCRejCorTen}{28.5}
\newcommand{\nTruncCorPi}{0.0108}
\newcommand{\nTruncCorGreen}{19}
\newcommand{\nGateStaticApplied}{259}
\newcommand{\nGateStaticSkipped}{53}
\newcommand{\nGateStaticAvoided}{17}
\newcommand{\nGateStaticForgone}{36}
\newcommand{\nGateStaticUpgrades}{181}
\newcommand{\nGateStaticUpgradesKept}{171}
\newcommand{\nGateStaticOracleApplied}{218}
\newcommand{\nGateStaticOracleSkipped}{94}
\newcommand{\nGateStaticOracleAvoided}{58}
\newcommand{\nGateStaticOracleForgone}{36}
\newcommand{\nGateStaticOracleUpgrades}{181}
\newcommand{\nGateStaticOracleUpgradesKept}{181}
\newcommand{\nDegradedStatic}{66}
\newcommand{\nDegradedStaticZoneUp}{6}
\newcommand{\nDegradedStaticNoChange}{60}
\newcommand{\nDegradedStaticAlreadyGreen}{59}
\newcommand{\nDegradedStaticSameNotGreen}{1}
\newcommand{\nDegradedStaticZoneDown}{0}
\newcommand{\nWellCalStaticN}{94}
\newcommand{\nWellCalStaticWorse}{58}
\newcommand{\nWellCalStaticMeanPct}{-1.3}
\newcommand{\nWellCalStaticWilcoxon}{8.9 \times 10^{-5}}
\newcommand{\nGateRollApplied}{259}
\newcommand{\nGateRollSkipped}{53}
\newcommand{\nGateRollAvoided}{44}
\newcommand{\nGateRollForgone}{9}
\newcommand{\nGateRollUpgrades}{204}
\newcommand{\nGateRollUpgradesKept}{184}
\newcommand{\nGateRollOracleApplied}{218}
\newcommand{\nGateRollOracleSkipped}{94}
\newcommand{\nGateRollOracleAvoided}{89}
\newcommand{\nGateRollOracleForgone}{5}
\newcommand{\nGateRollOracleUpgrades}{204}
\newcommand{\nGateRollOracleUpgradesKept}{204}
\newcommand{\nDegradedRoll}{174}
\newcommand{\nDegradedRollZoneUp}{71}
\newcommand{\nDegradedRollNoChange}{103}
\newcommand{\nDegradedRollAlreadyGreen}{100}
\newcommand{\nDegradedRollSameNotGreen}{1}
\newcommand{\nDegradedRollZoneDown}{2}
\newcommand{\nWellCalRollN}{94}
\newcommand{\nWellCalRollWorse}{89}
\newcommand{\nWellCalRollMeanPct}{-7.4}
\newcommand{\nWellCalRollWilcoxon}{10^{-16}}
\newcommand{\nGateRollAvoidedPct}{25.3}
\newcommand{\nGateRollOracleAvoidedPct}{51.1}
\newcommand{\nGateStaticAvoidedPct}{25.8}
\newcommand{\nGateRollUpgradesLost}{20}
\newcommand{\nGateStaticUpgradesLost}{10}
\newcommand{\nDegradedRollNoChangePct}{59.2}
\newcommand{\nDispDefaultSmall}{0.121}
\newcommand{\nDispDefaultMini}{0.115}
\newcommand{\nDispTwoHundredSmall}{0.331}
\newcommand{\nDispTwoHundredMini}{0.328}
\newcommand{\nDispThousandSmall}{0.958}
\newcommand{\nDispThousandMini}{0.945}
\newcommand{\nDispFullSmall}{1.087}
\newcommand{\nDispFullMini}{1.059}
\newcommand{\nDispTempLowSmall}{0.120}
\newcommand{\nDispTempLowMini}{0.114}
\newcommand{\nDispTempHighSmall}{0.121}
\newcommand{\nDispTempHighMini}{0.115}
\newcommand{\nDispNucleusLowSmall}{0.115}
\newcommand{\nDispNucleusLowMini}{0.109}
\newcommand{\nDoseCells}{1600}
\newcommand{\nPiSmallDefaultOne}{0.3884}
\newcommand{\nKupSmallDefaultOne}{0}
\newcommand{\nPiSmallDefaultTen}{0.4177}
\newcommand{\nKupSmallDefaultTen}{0}
\newcommand{\nRatioSmallDefaultOne}{38.837}
\newcommand{\nRatioSmallDefaultTen}{4.177}
\newcommand{\nPiSmallAnalyticOne}{0.0173}
\newcommand{\nKupSmallAnalyticOne}{8}
\newcommand{\nPiSmallAnalyticTen}{0.1027}
\newcommand{\nKupSmallAnalyticTen}{20}
\newcommand{\nRatioSmallAnalyticOne}{1.733}
\newcommand{\nRatioSmallAnalyticTen}{1.027}
\newcommand{\nPiMiniDefaultOne}{0.4188}
\newcommand{\nKupMiniDefaultOne}{0}
\newcommand{\nPiMiniDefaultTen}{0.4477}
\newcommand{\nKupMiniDefaultTen}{0}
\newcommand{\nRatioMiniDefaultOne}{41.884}
\newcommand{\nRatioMiniDefaultTen}{4.477}
\newcommand{\nPiMiniAnalyticOne}{0.0177}
\newcommand{\nKupMiniAnalyticOne}{8}
\newcommand{\nPiMiniAnalyticTen}{0.0981}
\newcommand{\nKupMiniAnalyticTen}{20}
\newcommand{\nRatioMiniAnalyticOne}{1.772}
\newcommand{\nRatioMiniAnalyticTen}{0.981}
\newcommand{\nGateBlocked}{4}
\newcommand{\nGateSeries}{13}
\newcommand{\nGateExtremesSmall}{19}
\newcommand{\nGateExtremesMini}{15}
\newcommand{\nWCDMPairs}{30}
\newcommand{\nWCDMSigBoot}{30}
\newcommand{\nWCDMSigAsymp}{29}
\newcommand{\nWCKupAsympGJR}{0.029}
\newcommand{\nWCKupBootGJR}{0.234}
\newcommand{\nLagForecasters}{13}
\newcommand{\nLagImmediate}{7}
\newcommand{\nLagMidDays}{77}
\newcommand{\nLagMidN}{4}
\newcommand{\nLagLateDays}{161}
\newcommand{\nLagLateN}{2}
\newcommand{\nLitClosureRMin}{0.005}
\newcommand{\nLitClosureRMax}{1.70}
\newcommand{\nGapDeltaMoment}{0.046}
\newcommand{\nGapUndGarchT}{10.6}
\newcommand{\nLitAnalyticDates}{40}
\newcommand{\nLitAnalyticSdPct}{0.55}
\newcommand{\nLitAnalyticGrid}{0.025}
\newcommand{\nRawPiGJRDefective}{0.0041}
\newcommand{\nRawPiGJRDefectivePct}{0.4}
\newcommand{\nPoolPiEWMA}{0.0119}
\newcommand{\nPoolPiMoiraiTwo}{0.0114}
\newcommand{\nPoolPiSmallAnalytic}{0.0113}
\newcommand{\nPoolPiGJR}{0.0111}
\newcommand{\nPoolBootEWMA}{0.058}
\newcommand{\nPoolClusterEWMA}{0.035}
\newcommand{\nuMCBShareWellSpec}{0.270}
\newcommand{\nuMCBShareAll}{0.329}
\newcommand{\nuMCBSharePct}{27.0}
\newcommand{\nLitTempGap}{0.001}
\newcommand{\nLitNucleusGap}{0.005}
\newcommand{\nLitSupportPct}{98.8}
\newcommand{\nLitRatioSmallTwoFive}{1.34}
\newcommand{\nLitRatioSmallFive}{1.14}
\newcommand{\nLitRatioMiniTwoFive}{1.32}
\newcommand{\nLitRatioMiniFive}{1.11}
\newcommand{\nLitAlphaResp}{1.08}
\newcommand{\nLitMainRMinAll}{0.001}
\newcommand{\nLitMainRMaxOK}{0.357}
\newcommand{\nLitMainROkN}{14}
\newcommand{\nLitQSMin}{4.67}
\newcommand{\nLitQSMax}{5.83}
\newcommand{\nLitWMin}{0.98}
\newcommand{\nLitWMax}{1.14}
\newcommand{\nLitCorPiTypical}{0.011}
\newcommand{\nLitMoiraiGapPP}{0.24}
\newcommand{\nLitLambdaTail}{1.9\times 10^{-7}}
\newcommand{\nLitQSDefaultSmall}{5.81}
\newcommand{\nLitQSAnalyticSmall}{5.03}
\newcommand{\nLitQSDefaultMini}{5.69}
\newcommand{\nLitQSAnalyticMini}{5.00}
\newcommand{\nSeqDQCells}{312}
\newcommand{\nSeqDQRejRaw}{81.7}
\newcommand{\nSeqDQRejCor}{53.8}
\newcommand{\nGapDeltaFree}{0.218}
\newcommand{\nGapUndFree}{75.5}
\newcommand{\nGapDeltaUni}{0.066}
\newcommand{\nGapUndUni}{49.4}
\newcommand{\nGapDeltaGateNow}{0.024}
\newcommand{\nGapUndGateNow}{30.9}
\newcommand{\nGapBandNow}{-1.800}
\newcommand{\nGapBandLower}{-3.500}
\newcommand{\nGapCellGoodWorst}{-1.947}
\newcommand{\nGapCellGoodBest}{-2.735}
\newcommand{\nGapCellTruncBest}{-0.283}
\newcommand{\nGapCellTruncWorst}{-0.074}
\newcommand{\nGapGapWidth}{1.664}
\newcommand{\nGapMarginUpper}{0.147}
\newcommand{\nGapMarginLower}{0.765}
\newcommand{\nGapCellsLowerBlocks}{0}
\newcommand{\nGapCellsTotal}{312}
\newcommand{\nGapCellsTrunc}{48}
\newcommand{\nGapDeltaChronos}{0.388}
\newcommand{\nGapQTrueStd}{-2.606}
\newcommand{\nGapUndTight}{25.6}
\newcommand{\nGapMarginTight}{0.007}
\newcommand{\nGapBandTight}{-1.940}
\newcommand{\nGapDMt}{0.399}
\newcommand{\nGapDMp}{0.69}
\newcommand{\nGapQSGapPct}{0.008}
\newcommand{\nGapVaRHonest}{2.61}
\newcommand{\nGapVaRAlt}{1.32}
\newcommand{\nGapRhoLo}{0.18}
\newcommand{\nGapRhoHi}{0.67}
\newcommand{\nGapDeltaHatLo}{0.109}
\newcommand{\nGapDeltaHatHi}{0.138}
\newcommand{\nGapEmpCoverage}{98.9}
\newcommand{\nGapGapAblCovid}{0.0058}
\newcommand{\nGapGapAblFull}{0.0005}
\newcommand{\nGapFisherKupiec}{0.49}
\newcommand{\nGapFisherSevere}{0.00078}
\newcommand{\nGapAblRhoLo}{+0.05}
\newcommand{\nGapAblRhoHi}{+0.46}
\newcommand{\nGapAblPairs}{4}
\newcommand{\nTLTau}{0.016}
\newcommand{\nTLTauOverAlpha}{1.6}
\newcommand{\nTLUndTFive}{12.4}
\newcommand{\nTLUndTThree}{16.3}
\newcommand{\nTLUndNormal}{7.8}
\newcommand{\nMCGreenTFiveSmall}{61.6}
\newcommand{\nMCGreenTFiveLarge}{74.2}
\newcommand{\nMCPiTFive}{0.0148}
\newcommand{\nMCGreenSkewSmall}{36.4}
\newcommand{\nMCGreenSkewLarge}{0.4}
\newcommand{\nMCTMax}{10{,}000}
\newcommand{\nMCTMin}{500}
\newcommand{\nMCReps}{500}
\newcommand{\nMCDgps}{5}
\newcommand{\nGateRollUpgradesLostWorse}{11}
\newcommand{\nGateRollUpgradesLostNet}{9}
\newcommand{\nGateRollUpgradeAndDeteriorate}{71}
\newcommand{\nSpearmanRPi}{0.991}
\newcommand{\nSpearmanRPiN}{14}


%─── Colours ─────────────────────────────────────────────────────────────────
\definecolor{pastelgreen}{rgb}{0.0,0.6,0.0}

%─── Revision markup (R1) ────────────────────────────────────────────────────
% Added or rewritten text is wrapped in \revised{...} and shown in blue.
% Revision colouring, off for the clean version. Set to the coloured
% variant only when a referee has asked to see what changed.
\newcommand{\revised}[1]{#1}
% \renewcommand{\revised}[1]{{\color{blue}#1}}
% Camera-ready toggle: uncomment the next line to render all revisions in black.
%\renewcommand{\revised}[1]{#1}

%─── Theorem environments ────────────────────────────────────────────────────
\newtheorem{theorem}{Theorem}[section]
\newtheorem{corollary}[theorem]{Corollary}
\newtheorem{definition}[theorem]{Definition}
\newtheorem{remark}[theorem]{Remark}
\newtheorem{assumption}[theorem]{Assumption}
\newtheorem{proposition}[theorem]{Proposition}
\newtheorem{lemma}[theorem]{Lemma}

%─── Notation ────────────────────────────────────────────────────────────────
\newcommand{\VaR}{\mathrm{VaR}}
\newcommand{\ES}{\mathrm{ES}}
\newcommand{\E}{\mathbb{E}}
\renewcommand{\P}{\mathsf{P}}
\newcommand{\1}{\boldsymbol{1}}
\newcommand{\qVstat}{\hat{q}_{\scriptscriptstyle\mathrm{V}}^{\,\mathrm{stat}}}
% Backward-compatible alias: three robustness fragments predate the
% \qV -> \qVstat rename and still emit \qV. Aliasing here avoids editing
% generated files, which would be overwritten on the next pipeline run.
\providecommand{\qV}{\qVstat}
\newcommand{\qVroll}{\hat{q}_{{\scriptscriptstyle\mathrm{V}},t}^{\,\mathrm{roll}}}
\newcommand{\qE}{\hat{q}_{\scriptscriptstyle\mathrm{E}}}
\newcommand{\EScqr}{\widehat{\ES}^{\,\mathrm{cqr}}}
\newcommand{\CI}{\hat{C}}
\newcommand{\Fhat}{\hat{\mathcal{F}}}
\newcommand{\quantlet}[1]{\href{#1}{\raisebox{-0.15em}{\includegraphics[height=0.9em]{ql_logo.png}}}}
\newcommand{\quantinar}[1]{\href{#1}{\raisebox{-0.15em}{\includegraphics[height=0.9em]{qr_logo.png}}}}

%══════════════════════════════════════════════════════════════════════════════
\begin{document}
	
	\begin{frontmatter}
		
		%─── PART 1: Title and Metadata ──────────────────────────────────────────────
\title{What Backtests Cannot Diagnose: Structural Validation of Tail-Risk Forecasting Pipelines}

		\author[1,2,3]{Daniel Traian Pele\corref{cor1}}
		\ead{danpele@ase.ro}
		\cortext[cor1]{Corresponding author}
		
		\author[1]{Vlad Bolov\u{a}neanu}
		\ead{vlad.bolovaneanu@ase.ro}
		
		\author[1,3]{Andrei Theodor Ginavar}
		\ead{ginavarandrei19@stud.ase.ro}
		
		
		\author[1,4]{Stefan Lessmann}
		\ead{stefan.lessmann@hu-berlin.de}
		
		\author[1,3,4,5,6]{Wolfgang Karl H{\"a}rdle}
		\ead{haerdle@hu-berlin.de}
		\affiliation[1]{organization={Bucharest University of Economic
				Studies (ASE)}, city={Bucharest}, country={Romania}}
		
		\affiliation[2]{organization={Institute for Economic Forecasting,
				Romanian Academy}, city={Bucharest}, country={Romania}}
		
		\affiliation[3]{organization={IDA Institute of Digital Assets},
			city={Bucharest}, country={Romania}}
		
		\affiliation[4]{organization={School of Business and Economics, Humboldt
				University of Berlin},
			city={Berlin}, country={Germany}}
		
		\affiliation[5]{organization={School of Informatics, University of Edinburgh},
			city={Edinburgh}, country={UK}}
		
		\affiliation[6]{organization={Department of Information Management and
				Finance, National Yang Ming Chiao Tung University}, city={Hsinchu},
			country={Taiwan}}                                         
		
		%─── PART 2: Abstract ───────────────────────────────────────────────────────
\begin{abstract}

	A tail-risk backtest reports \emph{that} a Value-at-Risk series miscovers. It cannot report \emph{why}. We show that this gap is where forecasting pipelines fail, and that the failures are neither rare nor exotic.

	We audit a \nMainAssets-asset, \nMainForecasters-forecaster panel and trace five defects in our own pipeline to specific lines of code: two sign inversions, an unstandardised Student-$t$ quantile in a GARCH benchmark, and two data-handling faults. The cleanest case is in published third-party code and reproduces from a public checkpoint without our data. Chronos-T5 represents its one-step predictive law as a categorical distribution over 4093 quantised bins; its packaged sampler defaults to $\texttt{top\_k} = 50$, discarding 98.8\% of that support, from the tails inward, before any quantile is read. The median is barely moved, so point-forecast evaluation is normal, while the 1\% violation rate reaches \nPiSmallDefaultOne{} and \nPiMiniDefaultOne. A controlled dose--response on identical contexts and seeds establishes the mechanism, and the parameters a practitioner would vary --- temperature, nucleus sampling --- move predictive dispersion by 0.001 and 0.006. Reading the law in closed form instead of sampling it has no sampling parameters, costs one forward pass, and moves the rate to \nPiSmallAnalyticOne{} and \nPiMiniAnalyticOne.

	Such failures survive the instruments applied to them. Kupiec's test rejects \nSeqKupiecRejRawOne\% of pairs at $\alpha = 0.01$, defective and merely imperfect series alike; the Christoffersen statistic is degenerate on \nSeqCCUndefRawOne\% of raw pairs and \nSeqCCUndefCorOne\% after recalibration, where its asymptotics do not hold; and a conformal shift targets marginal coverage whatever forecaster it corrects, so on a sign-inverted series it delivered 19 of \nMainAssets{} Basel green zones and every backtest downstream of it passed.

	We therefore validate structurally, before a series is calibrated or evaluated: two invariants that follow from what a quantile forecast is, one invariant tested through a proxy, and seven plausibility bands whose thresholds we state rather than derive. The gate blocks \nGateBlocked{} of \nGateSeries{} series here. Correctly configured, the zero-shot foundation models in this panel are neither catastrophically miscalibrated nor competitive --- they sit alongside GARCH-family benchmarks at the 1\% level, and a per-asset CAViaR fitted the ordinary way outperforms all of them.

	\bigskip\noindent
	\textbf{Keywords:} Value-at-Risk; Structural validation; Time series
	foundation models; Conformal prediction; Regulatory backtesting;
	Reproducibility

	\bigskip\noindent
	\textbf{JEL:} C14\,|\,C53\,|\,C58\,|\,G17\,|\,G32
\end{abstract}
\end{frontmatter}
	%══════════════════════════════════════════════════════════════════════════════
	%─── PART 4: Introduction ───────────────────────────────────────────────────
	%═════════════════════════════════════════
\section{Introduction}
\label{sec:intro}

A Value-at-Risk forecast is not a model output. It is the result of reducing a
model output --- a set of sampled paths, a grid of predictive quantiles, a fitted
parametric law, an order statistic over a window --- to a single number, and that
reduction is performed by code that the forecasting literature does not treat as
an object of study. This paper reports what happens when it is studied: in a
24-asset, 16-forecaster Value-at-Risk panel, the dominant driver of measured
extreme-tail failure is the configuration of that reduction step, not the
architecture of the model, its pretraining corpus, or the format in which it
exposes its predictive distribution.

The evidence is unusually direct, because the defects are our own. An audit of
the pipeline behind an earlier version of this paper traced five errors to
specific lines of code: two foundation-model series stored the $\alpha$-quantile
with its sign inverted; one classical GJR-GARCH benchmark applied an
unstandardised Student-$t$ quantile with its degrees of freedom hard-coded; and
two data-handling faults propagated a ticker alias and a stale forecast vintage.
Each is arithmetic rather than judgement, and each is documented with the
diagnostic that establishes it in the replication package. Their joint effect was
not noise. They moved reported 1\% violation rates by factors between 4 and 70,
and they moved them in both directions --- the GJR-GARCH defect inflated every
threshold by roughly 45\%, producing a 0.42\% violation rate against a nominal
1\% and placing the benchmark in the Basel green zone on all 24 assets. A defect
that flatters its owner attracts no scrutiny; that one was found last.

We develop the argument on the case that is not ours, because it is in published
third-party code and reproduces from a public checkpoint without our data.
Chronos-T5 \citep{ansari2024chronos} represents its predictive distribution as a
categorical law over 4093 quantised value bins. Its packaged sampling routine
defaults to $\texttt{top\_k} = 50$, a value inherited from text generation, which
discards 98.8\% of that support before any quantile is computed --- and discards
it from the tails inward. The median is barely affected, so a point-forecast
evaluation of this configuration is entirely normal. The 1\% tail is destroyed:
the empirical violation rate reaches 0.3884 and 0.4188 for the two checkpoint
sizes we audit, a thirty-nine-fold exceedance that would place every asset in the
Basel red zone. This is not a bug in Chronos. It is a documented default, correct
for the task it was designed for, silently wrong for this one.

Three properties make such failures durable rather than merely possible. First,
they are invisible to point-forecast evaluation, which is where forecasting
models are normally validated. Second, they are invisible to a responsible
sensitivity analysis: varying temperature across a fourfold range moves
predictive dispersion by 0.001 and nucleus sampling by 0.006, so an analyst who
sweeps the parameters that are documented and discussed will correctly report the
model as robust. Third --- and this is the part that concerns the tail-risk
literature specifically --- they are invisible to the backtests that regulate
this domain.

That third claim is measured rather than asserted. At $\alpha = 0.01$ the
independence test of \citet{christoffersen1998} is \emph{undefined} on 34.6\% of
our 312 model--asset pairs, because a 1\% tail generates too few exceedances to
populate a transition table, and on 53.5\% of them once the forecasts have been
recalibrated --- which is where every backtest in the earlier version of this
paper was computed, since a correction that spaces violations out disables the
test further. The share falls monotonically to 0.0\% at $\alpha = 0.10$ under
both readings. A degenerate table is an absence of evidence, but the natural
implementation returns ``no rejection'', which reads as a pass. Unconditional
coverage \citep{kupiec1995} is better defined and rejects
\nSeqKupiecRejRawOne\% of pairs at $\alpha = 0.01$ against
\nSeqKupiecRejRawTen\% at $\alpha = 0.10$; when a test rejects the majority
of defensible models, its rejection of any one of them carries little
information, and a genuinely defective series is not distinguishable from a
merely imperfect one. Both instruments are weakest exactly at the level at which
the Basel traffic-light system operates \citep{basel1996}.

Post-hoc recalibration does not rescue the situation; it conceals it. A
split-conformal shift restores marginal coverage regardless of the forecaster it
is applied to --- exactly and in finite samples under exchangeability
\citep{vovk2005algorithmic,lei2018distribution}, and asymptotically under the
$\beta$-mixing conditions of Theorem~\ref{thm:dependent_conformal_var} --- which
require a separation between calibration and evaluation that our contiguous split
does not provide, as Remark~\ref{rem:gap_vs_protocol} sets out. Neither guarantee
therefore covers the estimator exactly as we run it. The distinction matters for
the theory and not at all for the argument of this paragraph: on our data the
correction delivers coverage regardless. Applied to a sign-inverted
series it placed 19 of 24 assets in the Basel green zone. Every backtest reported
in the earlier version of this paper was computed downstream of recalibration, so
every one of them passed. The errors survived because the method worked.

\paragraph{Contributions} First, we establish the mechanism of the sampling
defect by controlled experiment rather than inference, and remove it. Holding
contexts, seeds and weights fixed and varying only the sampling configuration,
predictive dispersion recovers monotonically as the truncation is relaxed and one
thousand draws under the default contain exactly 50 distinct values on every
asset and both checkpoints. We then show that sampling is unnecessary at a
one-step horizon, where the predictive law is categorical and its quantiles are
available in closed form, and give an estimator that reads them directly from the
model's logits: it agrees with full-vocabulary sampling to within Monte Carlo
error, has no $\texttt{num\_samples}$, no $\texttt{top\_k}$ and no seed, and
costs one forward pass. Under it the violation rate falls from 0.3884 and 0.4188
to 0.0175 and 0.0178.

Second, we quantify the blindness of the standard tail diagnostics, as above,
and separate two claims that are usually merged: that these tests are being
misapplied, and that they are being applied where they have the least power. It
is the second. This is the argument for validating a forecast series
structurally, before it is calibrated or evaluated.

Third, we propose such a validation gate --- ten checks on a forecast series,
of which two are invariants, one an invariant tested through a proxy, and seven
plausibility bands with thresholds we state rather than derive. Eight need
nothing but the series and the returns; two need an evaluation window, and we say
so rather than claim the gate is entirely free of one. We report its scoping
honestly in the other direction too: two of the checks are inapplicable to order-statistic estimators and
say so rather than returning a verdict, and all ten were written after the
defects were known, so their performance on those defects is in-sample by
construction. The gate blocks \nGateBlocked{} of \nGateSeries{} series in this panel.

Fourth, for the recalibration layer we retain an asymptotic coverage guarantee
for the single-split estimator under $\beta$-mixing dependence with explicit
gap-separation conditions, together with finite-window bounds for the rolling
variant under distributional drift, and we report the intervention's
\emph{indication} rather than its average effect. Recalibration is close to free
where it is not needed and harmful where the base forecast is already calibrated:
under the rolling estimator, \nWellCalRollWorse{} of the \nWellCalRollN{} pairs
that pass Kupiec on the test window are degraded by correcting them --- a
description of who is harmed, not yet a rule for deciding. Gating the correction on a failed raw backtest
helps, but far less than an in-sample evaluation of that rule suggests: keyed on
the calibration window, which is the only information available at the decision
date, it avoids \nGateRollAvoidedPct\% of the damage at $\alpha = 0.01$ and costs
\nGateRollUpgradesLost{} of \nGateRollUpgrades{} Basel zone upgrades. Keyed on
the evaluation window it would avoid \nGateRollOracleAvoidedPct\% and cost none,
and we report that only as the oracle it is.

\paragraph{What this paper does not claim} It does not claim that foundation
models are unsuitable for tail-risk forecasting, nor that they are ready for it.
Correctly configured, the six zero-shot foundation-model series in our panel are
neither catastrophically miscalibrated nor competitive: their raw 1\% violation
rates run from 0.0143 to 0.0294, which places them alongside the GARCH-family
benchmarks on the same assets and windows, and a per-asset CAViaR
\citep{engle2004caviar} fitted the ordinary way outperforms all of them. It does
not claim that configuration failure explains tail miscalibration in general ---
we observe five defects in one pipeline, and the frequency of such defects in the
literature is unknown, though we note that neither of the two recent tail-risk
papers we surveyed that use sample-based foundation models reports a single
generation parameter. And it does not claim that the ten structural checks are
sufficient, and only three of the ten are necessary conditions in the strict
sense: the other seven are plausibility bands. What can be said for all of them
is that they are cheap, that each states a property a practitioner would endorse
before seeing any data, and that the alternative --- coverage --- cannot do this
job at all.

The remainder of the paper is organised as follows.
Section~\ref{sec:lit} situates the contribution in the literature on tail
miscalibration, regulatory backtesting, and post-hoc calibration.
Section~\ref{sec:framework} develops the two-layer audit, with asymptotic
coverage results under $\beta$-mixing dependence, and describes the asset
universe, the forecasters, and the experimental design.
Section~\ref{sec:mechanism} establishes the sampling mechanism by controlled
experiment and gives the analytic estimator that removes it.
Section~\ref{sec:results} presents the empirical audit on the corrected panel,
the ordering by $\bar R$, and the static-versus-rolling comparison.
Section~\ref{sec:blindspot} measures what the standard diagnostics can and cannot
detect, and Section~\ref{sec:gate} states the structural validation gate.
Section~\ref{sec:discuss} discusses calibration--sharpness trade-offs,
deployment considerations, and the FRTB capital arithmetic.
Section~\ref{sec:conclude} concludes. Proofs and supplementary tables
are in the appendices.

Data and code accompanying this paper are available via the Quantlet platform \quantlet{https://github.com/QuantLet/Conformal_Oracle/}; slides are on the Quantinar platform \quantinar{https://quantinar.com/course/100048/recalibrating-tail-risk-forecasts-under-temporal-dependence}.
	
%══════════════════════════════════════════════════════════════════════════════
%─── PART 5: Literature Review ──────────────────────────────────────────────
%══════════════════════════════════════════════════════════════════════════════
\section{Related Literature}
\label{sec:lit}

This paper relates to three strands: tail miscalibration in probabilistic forecasters, regulatory backtesting of tail-risk forecasts, and post-hoc calibration under temporal dependence.

Tail miscalibration is well-documented across forecaster classes. Parametric volatility models can systematically under- or over-estimate extreme quantiles under heavy-tailed innovations or regime shifts \citep{christoffersen1998,mcneil2000estimation,nolde2017elicitability}. Recent time series foundation models (TSFMs) --- including Chronos \citep{ansari2024chronos,ansari2025chronos2}, TimesFM \citep{das2024timesfm}, Moirai \citep{woo2024moirai,liu2025moirai}, and Lag-Llama \citep{rasul2024lagllama} --- achieve competitive zero-shot forecasting performance but exhibit a similar pattern: \citet{goel2024tsfmvar} show that pre-trained quantiles require fine-tuning for acceptable Value-at-Risk (VaR) coverage, and \citet{pele2025llmvar} identify tail compression in instruction-tuned models, with similar miscalibration documented for volatility forecasting \citep{goel2025volatility,rahimikia2025revisiting}. These findings, spanning parametric and foundation-model regimes alike, motivate post-hoc calibration but do not address the identification constraint imposed by extreme-tail sparsity, and they stop short of a taxonomy of when post-hoc correction adjusts an informative base versus when it acts as an empirical fallback.

Regulatory backtesting procedures \citep{kupiec1995,christoffersen1998,basel1996,acerbi2014back,fissler2016higher,nolde2017elicitability} detect violations of nominal coverage, exceedance dependence, and tail-loss inconsistency, but do not prescribe forecast corrections.

Post-hoc calibration adjusts the output of a fixed forecaster without retraining \citep{gneiting2013combining,taillardat2016calibrated}. These procedures sit within the broader probabilistic-forecasting literature on calibration--sharpness trade-offs \citep{gneiting2007probabilistic,gneiting2014probabilistic}, where calibration restoration is balanced against sharpness loss. Model-based alternatives, such as CAViaR \citep{engle2004caviar} and conditional-density frameworks such as the Generalised Additive Model for Location, Scale and Shape \citep[GAMLSS;][]{rigby2005generalized}, estimate tail behaviour directly but impose structure on the conditional distribution. Conformal prediction \citep{vovk2005algorithmic}, applied to quantile regression by \citet{romano2019conformalized} as conformalized quantile regression (CQR), provides a model-free alternative. Standard guarantees rely on exchangeability, which is violated in financial time series by volatility clustering, serial dependence, and structural change; extensions to dependent or non-exchangeable data include weighted conformal methods \citep{barber2023conformal}, blocking schemes \citep{chernozhukov2018exact}, and adaptive procedures \citep{zaffran2022conformal,xu2024conformal,angelopoulos2024pid}.

The TSFM literature receives closer attention above because this class is recent and has not yet been the subject of a systematic 1\% tail audit; the parametric volatility literature is mature and enters this paper through standard benchmarks rather than fresh review.

Building on these strands, this paper differs from all of them in where it looks. \citet{goel2024tsfmvar} and \citet{goel2025volatility} reduce TSFM tail miscalibration by fine-tuning; \citet{gibbs2021adaptive} deliver adaptive coverage under nonstationarity by online updates of the miscoverage rate; \citet{tamingtailrisk2026} extend conformal risk control to nonstationary portfolio VaR through regime-similarity weighting. Each treats the forecast series as given and improves what is done to it. We ask instead whether the series is what it is claimed to be, and find that in our own pipeline it repeatedly was not --- for reasons that lie in the code that converts a predictive distribution into a threshold, and that no procedure downstream of that conversion can detect.

There is a further gap this paper can only point at. Neither of the two recent tail-risk studies we surveyed that use sample-based foundation models reports a single generation parameter --- no \texttt{num\_samples}, \texttt{top\_k}, \texttt{top\_p}, temperature or seed appears in either. \citet{rahimikia2025revisiting} is the clean case: it is centrally about foundation models in finance, Chronos is its most-used model, and none of these parameters occurs anywhere in the text. Two papers do not establish a practice, and we do not claim they do; what they establish is that the reporting convention which would make this class of defect visible to a reader does not currently exist.

	%══════════════════════════════════════════════════════════════════════════════
	%─── PART 6: Framework ──────────────────────────────────────────────────────
	%══════════════════════════════════════════════════════════════════════════════
	
\section{Methodology}
\label{sec:framework}

The audit has two layers: a base model $M$ producing one-step-ahead probabilistic forecasts (Section~\ref{sec:layer1}), and a one-parameter conformal shift that recalibrates the predicted lower quantile (Section~\ref{sec:layer2}). Two implementations of the second layer are developed in Section~\ref{sec:layer2}: a single-split estimator, which admits asymptotic coverage under $\beta$-mixing dependence, and a rolling variant.

\paragraph*{Notation}
Let $r_t$ denote the daily log return and $\hat q_t^{lo} = (\widehat F_t^M)^{-1}(\alpha)$ the predicted lower $\alpha$-quantile from base model~$M$; we fix $\alpha = 0.01$. VaR is reported as a positive loss threshold:
\begin{equation}
	\label{eq:var_raw}
	\VaR_t^{\mathrm{raw}} = -\hat q_t^{lo}.
\end{equation}
The nonconformity score is
\begin{equation}
	\label{eq:score}
	S_t = \hat q_t^{lo} - r_t,
\end{equation}
and the conformal shift is denoted $q_V$, with single-split and rolling variants
\begin{equation}
	\label{eq:var_cp_roll}
	\VaR_t^{\mathrm{cp}} = \VaR_t^{\mathrm{raw}} + q_V^{\mathrm{stat}},
	\qquad
	\VaR_t^{\mathrm{roll}} = \VaR_t^{\mathrm{raw}} + q_V^{\mathrm{roll}}(t).
\end{equation}

\subsection{Layer~1: Probabilistic Forecasting}
\label{sec:layer1}

\subsubsection{Base Model Setup}
\label{sec:setup}

The first layer produces one-step-ahead probabilistic forecasts of returns from a fixed base model~$M$.

For sample-based forecasters (Chronos-Small, Chronos-Mini, Moirai~1.1, and Lag-Llama), the forecast at time~$t$ consists of $N = 1{,}000$ Monte Carlo draws $\{\hat r_t^{i;M}\}_{i=1}^N$, yielding the empirical predictive CDF
\begin{equation}
	\label{eq:ecdf_sample}
	\widehat F_t^M(x)
	= \frac{1}{N}\sum_{i=1}^{N}
	\mathbf{1}\{\hat r_t^{i;M}\leq x\}.
\end{equation}
The raw VaR forecast is the empirical lower $\alpha$-quantile,
\begin{equation}
	\label{eq:var_sample}
	\widehat{\VaR}_t^{\alpha;M}
	= -(\widehat F_t^M)^{-1}(\alpha).
\end{equation}

For quantile-grid forecasters (TimesFM~2.5 and Moirai~2.0), the predictive object is a finite grid of conditional quantiles $\{\hat q_t(u_j)\}_{j=1}^{J}$, with $J = 9$ deciles. Since the audit operates at the 1\% lower-tail level, we obtain $\hat q_t(0.01)$ by fitting a three-parameter Student-$t$ distribution to the native predicted quantiles:
\begin{equation}
	\label{eq:tail_completion}
	(\hat\nu_t, \hat\mu_t, \hat\sigma_t)
	=
	\arg\min_{\nu, \mu, \sigma}
	\sum_{j=1}^{J}
	\left[
	\hat q_t(u_j)
	-
	\left\{
	\mu + \sigma\, t_{\nu}^{-1}(u_j)
	\right\}
	\right]^2,
	\qquad \nu > 2,\ \sigma > 0,
\end{equation}
and set
\begin{equation}
	\label{eq:var_grid}
	\hat q_t(0.01)
	=
	\hat\mu_t + \hat\sigma_t\, t_{\hat\nu_t}^{-1}(0.01).
\end{equation}
The Student-$t$ closure is a parsimonious heavy-tailed default in financial VaR modelling \citep{bollerslev1987,hansen1994autoregressive,mcneil2000estimation}. It is not, however, an innocuous choice: Supplement~S.10 shows that the closure rule governs the tail behaviour of a grid model far more than the model does, so this decision is reported as part of the forecaster rather than as an implementation detail.

For classical benchmarks (GJR-GARCH, GARCH-Normal, EWMA, Historical Simulation), VaR is computed from the model-implied conditional distribution or, for Historical Simulation, the empirical 250-day rolling distribution.

Expected Shortfall is reported separately at $\alpha = 0.025$ following the Basel III/FRTB convention \citep{bcbs2019frtb} and is computed from the same predictive object as VaR. Because ES is not elicitable as a standalone functional \citep{gneiting2011making}, the conformal construction is developed for VaR only; ES computation details and the heuristic ES shift are in Supplement~S.1.

\subsubsection{Diagnosing Miscalibration}
\label{sec:diag}

We apply the standard regulatory backtests --- Kupiec~\citep{kupiec1995} for unconditional coverage, Christoffersen~\citep{christoffersen1998} for joint coverage and serial independence, and the Basel Traffic Light \citep{basel1996} for zone classification --- and define the violation indicator
\begin{equation}
	V_t = \mathbf{1}\bigl(r_t < -\widehat{\VaR}_t^{\alpha;M}\bigr).
\end{equation}

Tail-loss consistency is assessed through the $Z_2$ statistic \citep{acerbi2014back}, computed with the time-averaged ES denominator $Z_2^{\mathrm{mod}}$, which avoids division by near-zero daily ES on quiet days (Supplement~S.3). Our implementation of it carried a sign defect of its own, found late and reported in Supplement~S.1 rather than quietly repaired. Beyond coverage, the Fissler--Ziegel scoring rule \citep{fissler2016higher} ranks $(\VaR, \ES)$ jointly; full construction is in Supplement~S.1. The conformal shift $\qVstat$ introduced in Section~\ref{sec:layer2} complements these tests by quantifying miscalibration on a continuous, signed scale; the ratio $\bar R$ built from it in Section~\ref{sec:R_meaning} is non-negative by construction and carries magnitude only.

Beyond the per-asset backtests, four panel-level inferential procedures support the empirical audit: a diagnostic regression of the Quantile Score improvement on the conformal threshold, a cross-sectional association between the threshold and asset characteristics, a Diebold--Mariano pairwise forecast comparison with Driscoll--Kraay panel-HAC variance, and a panel-pooled Kupiec test with cluster-robust inference. Specifications are in Supplement~S.3; results enter the empirical sections at first use.

\subsection{Layer~2: Conformal Recalibration}
\label{sec:layer2}

\subsubsection{One-Sided Conformal VaR Correction}
\label{sec:cqr_var}

Conformal prediction provides a model-free framework for post-hoc calibration of predictive distributions \citep{vovk2005algorithmic,romano2019conformalized}. We adapt it to one-sided tail risk: forecast outputs are split chronologically into calibration and test sets with fraction $f_c = 0.70$ (robustness in Supplement~S.4).

The nonconformity score $S_t$ from equation~\eqref{eq:score} is positive when the realised loss exceeds the predicted quantile. The conformal shift is the empirical $(1-\alpha)$-quantile of the calibration scores,
\begin{equation}
	\label{eq:qV}
	\qVstat
	= Q_{1-\alpha}\!\left(\{S_t\}_{t\in\mathcal D_n}\right),
\end{equation}
where $\mathcal D_n$ is the calibration sample. The recalibrated lower quantile and VaR are then given by equation~\eqref{eq:var_cp_roll} with $q_V^{\mathrm{stat}} = \qVstat$. A positive $\qVstat$ shifts the predicted quantile downward, increasing reported VaR and correcting underestimation of tail risk; a negative value indicates over-conservatism.

The construction makes no parametric assumption on returns; the coverage guarantee under temporal dependence is developed in Section~\ref{sec:validity}.

The one-sided design is dictated by the regulatory loss function, not by a symmetry consideration. Basel III/FRTB capital and the Yellow/Red zone triggers are functions of loss-side exceedances alone (Section~\ref{sec:economic}): a day on which the realised return falls below the reported VaR is counted, a day on which it comfortably exceeds it is not. A correction applied to the upper tail would therefore alter no capital charge and no zone assignment, while widening the reported interval and degrading sharpness. We accordingly shift only $\hat q_t^{lo}$ and leave the remainder of the predictive distribution untouched; two-sided and distributional alternatives are evaluated in Section~\ref{sec:baselines}. Operational safeguards, including the non-negativity constraint $\qVstat \geq 0$, are discussed in Section~\ref{sec:failure}.

\begin{remark}[Variance scaling under tail sparsity]
	\label{rem:tail_sparsity}
	Tail calibration relies on events $\{r_t \leq \hat q_t^{lo}\}$, whose expected number is $\E[N_\alpha] = \alpha T$. Estimators fitted to this tail subset inherit variance of order $k/N_\alpha$ for $k$ parameters. By contrast, $\qVstat$ is a single order statistic of the full calibration sample. A one-parameter correction therefore enjoys a favourable bias--variance trade-off under tail sparsity, and the correction magnitude is interpretable as a diagnostic precisely because the construction adds no further parameters. Values of $\qVstat$ that dominate the scale of $\hat q_t^{lo}$ signal the high-$\bar R$ range.
\end{remark}

\subsubsection{Validity under Temporal Dependence}
\label{sec:validity}

Standard conformal coverage guarantees require exchangeability, which volatility clustering and serial dependence in financial returns violate. We therefore state an asymptotic coverage result under weak dependence.\footnote{We work with $\beta$-mixing rather than $\alpha$-mixing because the conformal threshold is an empirical quantile of the score process and the proof uses blocking/coupling arguments controlling total-variation distance between calibration and evaluation blocks.}

\begin{assumption}\label{ass:main}
	Let $S_t = \hat q_t^{lo} - r_t$ denote the nonconformity score process and let $n = \lfloor f_c T \rfloor$ be the calibration sample size.
	\begin{enumerate}[label=\textup{(A\arabic*)}]
		\item\label{ass:stationary}
		$\{S_t\}$ is strictly stationary and geometrically $\beta$-mixing: $\beta(k) \le C\rho^k$, $\rho \in (0,1)$.
		\item\label{ass:density}
		The marginal distribution $F_S$ admits a density $f_S$ that is continuous and bounded away from zero near the $(1-\alpha)$-quantile $q_{1-\alpha}$.
		\item\label{ass:weak_dep}
		There exists a gap sequence $g_n \to \infty$, $g_n/n \to 0$, such that the test score $S_{n+g_n+1}$ is evaluated after the calibration sample $\mathcal D_n = (S_1, \ldots, S_n)$, separated by $g_n$ unused observations.
	\end{enumerate}
\end{assumption}

\begin{theorem}
	\label{thm:dependent_conformal_var}
	Under Assumption~\ref{ass:main}, the single-split conformal VaR estimator satisfies
	\begin{equation}
		\label{eq:coverage_bound}
		\mathbb P\!\left(
		r_{n+g_n+1} \ge \hat q_{n+g_n+1}^{lo} - \qVstat
		\right)
		\ge
		1 - \alpha - \Delta_n,
	\end{equation}
	where $\Delta_n \le C\sqrt{(\log n)/n} + \beta(g_n)$ and $C > 0$ depends on the density near $q_{1-\alpha}$ and on the mixing rate.
\end{theorem}

\begin{remark}[What the theorem assumes, and what the protocol does]
	\label{rem:gap_vs_protocol}
	Two gaps between this result and the empirical design are worth stating
	rather than leaving for a reader to find.

	First, the theorem requires a separation $g_n\to\infty$ between the
	calibration block and the evaluation point, and by
	Lemma~S.9.1 that separation has to exceed the forecaster's
	context length $m$ before the mixing bound bites at all. The protocol used
	throughout the empirical sections is a \emph{contiguous} $70/30$ split:
	$g_n = 0$. The theorem therefore does not apply to the estimator as run. We
	report the gap ablation for exactly this reason (Supplement~S.4): inserting
	a gap moves $\hat\pi$ by at most $0.0005$ over the full sample and $0.0058$
	within the COVID window, so the cost of the discrepancy is measurable and
	small on this panel --- which is evidence, not a proof, and is offered as
	such.

	Second, the finite-sample conformal guarantee is a statement about
	exchangeable scores. Financial return series are not exchangeable, which is
	what this theorem exists to address; its conclusion is asymptotic, with an
	explicit remainder, and no finite-sample coverage claim is made anywhere in
	this paper for the dependent case.
\end{remark}

The bound is conservative at typical financial sample sizes. On six representative model--asset pairs spanning estimated score persistence $\hat\rho \in [0.18, 0.67]$, the estimated remainder $\hat\Delta_n$ lies between 0.109 and 0.138 --- a coverage floor of approximately 88\% against empirical coverage averaging 98.9\%. Across all 312 model--asset pairs, empirical coverage exceeds the corresponding lower bound $1 - \alpha - \hat\Delta_n$: the worst corrected coverage in the panel is $0.968$, against a floor of at most $0.891$. The bound provides a dependence-robust justification for the single-split correction; finite-sample performance is assessed in Section~\ref{sec:results}.

The remainder simplifies under standard GARCH dynamics.

\begin{corollary}
	\label{cor:garch_conformal_var}
	Let $\{r_t\}$ follow a GARCH(1,1) process satisfying $\E[\log(\alpha_1\varepsilon_t^2 + \beta_1)] < 0$. Then the score process $\{S_t\}$ is strictly stationary and geometrically $\beta$-mixing with $\beta(k) \le a\rho^k$. Choosing $g_n = c\log n$ with $c > 1/|\log\rho|$, Theorem~\ref{thm:dependent_conformal_var} applies with $\Delta_n = \mathcal O(\sqrt{(\log n)/n})$.
\end{corollary}


\subsubsection{Rolling Conformal Calibration}
\label{sec:rolling_method}

Financial time series exhibit time-varying volatility and structural change, motivating a locally adaptive conformal correction. In the rolling implementation, the conformal threshold is recomputed at each time $t$ from a trailing window of size $w = 250$:
\begin{equation}
	\label{eq:qV_roll}
	\qVroll
	=
	Q_{1-\alpha}\!\left(\{S_\tau\}_{\tau=t-w}^{t-1}\right),
\end{equation}
with the recalibrated lower quantile and VaR given by equation~\eqref{eq:var_cp_roll} substituting $q_V^{\mathrm{roll}}(t) = \qVroll$.

The rolling variant is an operational heuristic, not a construction with finite-sample conformal validity: Theorem~\ref{thm:dependent_conformal_var} covers the single-split estimator only, and recomputing the threshold from a trailing window violates the single-split assumption.\footnote{Complementary results from the online learning literature \citep{zinkevich2003online} and the weighted conformal framework \citep{barber2023conformal,angelopoulos2024pid} bound time-average miscoverage and finite-window deviation under drift, but cover related objects --- exact OGD updates and exponentially weighted quantiles, respectively --- not the uniform-weight rolling quantile used here. PropositionsProposition~S.9.4 andProposition~S.9.6 in Supplement~S.9 adapt these bounds; the proofs are included for completeness.} Empirical support comes from the simulation study in Section~\ref{sec:simulation} and from the drift diagnostic $\hat\delta_w(t)$ (Supplement~S.9.5), which flags windows where local recalibration departs from the global one.

The formal conformal guarantee in Theorem~\ref{thm:dependent_conformal_var} applies to the single-split estimator. The rolling estimator should be read as a locally adaptive engineering variant whose validity is assessed empirically and supported by the time-average and drift bounds in Supplement~S.9, not as an exact conformal procedure.

\subsubsection{Alternative Recalibration Baselines}
\label{sec:alternatives}

We compare the one-parameter conformal shift against alternatives that vary on three axes: parameter count, correction structure (location only, location--scale, monotone, fully nonparametric), and target object (correction, conditional distribution, or tail). We compare against nine baselines: five single-shot post-hoc corrections, two distributional or nonlinear methods (GAMLSS-SST and gradient-boosted quantile regression), and two dedicated VaR models (EVT-POT and filtered historical simulation). Full specifications are in Supplement~S.2; Table~S.2 of the supplement provides a compact reference. The comparison tests Remark~\ref{rem:tail_sparsity}: at $\alpha = 0.01$, additional parameters increase variance faster than they reduce bias. Empirical results are in Section~\ref{sec:baselines}.
	%══════════════════════════════════════════════════════════════════════════════
	%─── PART 7: Data and Experimental Setup ────────────────────────────────────
	%══════════════════════════════════════════════════════════════════════════════
\subsection{Data and Models}
\label{sec:data_models}

\subsubsection{Data}
\label{sec:data}

The empirical analysis covers 24 financial assets across five classes: 11 equity indices, 3 government bond exchange-traded funds (ETFs), 4 commodities, 2 cryptocurrencies, and 4 foreign exchange (FX) pairs. Daily log-returns are computed from adjusted closing prices obtained from Yahoo Finance over January 2000--March 2026. Price series were retrieved on 13 March 2026. Cryptocurrency histories begin in 2014 for Bitcoin and 2017 for Ethereum.

Returns are computed on each asset's native trading calendar without cross-asset alignment, and no further data cleaning is applied: missing days are dropped rather than forward-filled, and asynchronous closes across classes are retained. The evaluation therefore reflects deployment conditions rather than a synchronised research panel.

Each asset's return series is partitioned chronologically: the single-split estimator uses the first 70\% as calibration and the remaining 30\% as the out-of-sample test window ($f_c = 0.70$, Section~\ref{sec:cqr_var}); the rolling estimator recomputes the threshold at each $t$ from a trailing window of size $w = 250$ (Section~\ref{sec:rolling_method}) and is evaluated on the same 30\% test window. The full asset list, per-asset sample sizes, and per-asset test windows are in Supplement~S.7.

\subsubsection{Models}
\label{sec:models}

We evaluate six TSFM forecasters in strict zero-shot mode (no updating or fine-tuning) against four classical benchmarks estimated per asset. Fine-tuning TSFMs is left to future work.

The TSFM set spans two predictive interfaces. Sample-based forecasters (Chronos-Small, Chronos-Mini \citep{ansari2024chronos,ansari2025chronos2} with tokenised sequence models; Lag-Llama \citep{rasul2024lagllama} with a LLaMA-inspired architecture and Student-$t$ heads; Moirai~1.1 \citep{woo2024moirai}) return $N = 1{,}000$ Monte Carlo draws from the predictive distribution. Quantile-grid forecasters (TimesFM~2.5 \citep{das2024timesfm}; Moirai~2.0 \citep{liu2025moirai}) return a fixed nine-decile grid, from which a 1\% quantile must be supplied by a tail-completion step. Each Chronos checkpoint enters the panel \emph{twice}: once as the shipped series, sampled under the checkpoint's own default configuration, and once read analytically from the same weights by the estimator of Section~\ref{sec:mechanism:analytic}. The pair is the paper's controlled comparison, and it holds the model fixed while varying only the code that turns its output into a threshold.

The classical benchmarks --- GJR-GARCH \citep{glosten1993relation}, GARCH-Normal \citep{bollerslev1986generalized}, EWMA \citep{riskmetrics1996}, and Historical Simulation --- represent standard volatility and tail-risk models; estimation details are in Supplement~S.8 and Section~\ref{sec:exp_design}.

\subsubsection{Experimental Design}
\label{sec:exp_design}

\textbf{Two panels, named once.} The paper reports two, and mixing them is the
easiest way to produce a number nobody can reconcile, so they are defined here
and tagged wherever they appear.

\emph{The main panel} is \nMainForecasters{} forecasters $\times$ \nMainAssets{}
assets $=$ \nMainPairs{} pairs at $\alpha = 0.01$. This is Table~\ref{tab:master}.
It includes CAViaR-AS, CAViaR-SAV and GAS-$t$, which are estimated directly from
returns by a separate routine and produce no stored forecast series.

\emph{The sequence panel} is \nSeqForecasters{} forecasters $\times$
\nMainAssets{} assets $=$ \nSeqPairs{} pairs, evaluated at four levels. It is the
subset with stored per-date series, and every analysis that needs one runs on it:
the Christoffersen degeneracy rates by level, the Diebold--Mariano comparisons,
the wild-cluster bootstrap, the panel-pooled coverage tests, the deterioration
counts behind the indication rule, and Figure~\ref{fig:traffic_light}.

The two agree wherever they overlap --- the sequence panel's static Basel zones
reproduce Table~\ref{tab:master} on all \nSeqPairs{} shared pairs --- and every
figure in the text is generated from the artefacts by
\texttt{scripts/paper\_numbers.py}, which fails the build if a value in the
manuscript drifts from the value in the data.

\textbf{The comparison is deliberately asymmetric, and we state that here rather than as a qualification later.} The classical and dynamic benchmarks are estimated per asset and re-estimated at every step; the foundation models are deployed strictly zero-shot, with no fine-tuning and no per-asset adaptation. These are not equivalent adaptation regimes, and any Quantile Score comparison between the two groups reflects that difference as well as any difference in model quality. We retain the asymmetry because it is the operationally relevant contrast --- a zero-shot model is what a practitioner can deploy without an estimation pipeline per instrument --- but no conclusion in this paper rests on a TSFM-versus-GARCH score ranking. The audit's findings are within-group: the interface contrast is between two releases of the same family under identical deployment, and the recalibration results are computed separately for every forecaster.

All forecasts are one-step-ahead. TSFMs use a fixed context length of 512 observations, the default for all six pretrained checkpoints (Supplement~S.8). Classical benchmarks are estimated on rolling windows of 250 observations, following the Basel convention \citep{basel1996}, with parameters re-estimated at each step.

A warm-up period (512 observations for TSFMs, 250 for classical benchmarks) precedes the calibration--test split. For full-history series (2000--2026), the test sample is approximately 1{,}500 observations per asset, yielding roughly 36{,}600 panel observations per model and 366 expected violations at $\alpha = 0.01$.
%══════════════════════════════════════════════════════════════════════════════
%─── PART 6b: The sampling mechanism ────────────────────────────────────────
%══════════════════════════════════════════════════════════════════════════════
\section{The Sampling Mechanism}
\label{sec:mechanism}

This section establishes, by controlled experiment, that a single documented
default in a checkpoint's packaged inference routine accounts for a
thirty-nine-fold exceedance of a nominal 1\% Value-at-Risk, and then removes the
sampler from the estimation path entirely. The case is developed on Chronos-T5
\citep{ansari2024chronos} because it is the one defect in this paper that lives
in published third-party code and reproduces from a public checkpoint without our
data. Everything in this section is therefore verifiable by a reader who has
neither our forecasts nor our returns.

\subsection{What the default does}
\label{sec:mechanism:default}

Chronos-T5 does not emit a parametric predictive law. It quantises the scaled
context into 4093 value bins and treats forecasting as next-token prediction over
that vocabulary, so its one-step predictive distribution is categorical with 4093
atoms. Predictions are drawn from it by the same sampling machinery used for
text, and that machinery has a $\texttt{top\_k}$ parameter, defaulted to 50:
before a single path is drawn, the support is restricted to the 50 most probable
bins and the remaining 4043 are assigned probability zero.

Two properties of this restriction matter here. It removes 98.8\% of the support,
and it removes it \emph{from the tails inward}, because the discarded bins are by
construction the least probable ones. A quantile read from the surviving support
is therefore a quantile of a different, radically narrower distribution --- while
the centre of that distribution, and hence any point forecast, is barely moved.

\subsection{Design}
\label{sec:mechanism:design}

The experiment varies one thing. For each of two checkpoints
(\texttt{chronos-t5-small}, \texttt{chronos-t5-mini}), four assets of unlike
character (SP500, GOLD, BTC, EURUSD) and 200 evenly spaced dates with a full
512-observation context, we draw 1000 samples at horizon one under eight sampling
configurations: $\texttt{top\_k} \in \{50, 200, 1000, 4094\}$,\footnote{The
vocabulary holds 4096 tokens, of which two are special, leaving 4093 value bins
plus the pad token. Setting $\texttt{top\_k} = 4094$ therefore disables the
truncation with one place to spare; it is not an off-by-one against the 4093 bins
quoted elsewhere in this paper.} then
$\texttt{top\_p} \in \{0.9, 0.99\}$ and temperature $\in \{0.5, 2.0\}$ with
$\texttt{top\_k}$ held at its default. Context, seed and weights are identical
across configurations, so the only thing that differs between cells is the
sampling rule.

Three readings were fixed before the experiment was run. If the truncation is the
mechanism, then (i) predictive dispersion recovers monotonically in
$\texttt{top\_k}$ and reaches its target only at full vocabulary, (ii) dispersion
is unmoved by temperature and by nucleus sampling, and (iii) the number of
distinct values among 1000 draws under the default equals $\texttt{top\_k}$
exactly, on every asset and both checkpoints.

\subsection{Result}
\label{sec:mechanism:result}

\begin{table}[htbp]
	\centering
	\caption{Sampling dose--response for Chronos-T5. Dispersion is the ratio of
		predictive standard deviation to the asset's 250-day realised volatility,
		averaged over the 800 (asset, date) cells of each configuration; 1.0 is
		the target. Context, seed and weights are identical across rows.}
	\label{tab:dose_response}
	% Generated by Quantlets/CO_chronos_sampling/run_dose_table.py
\begin{tabular}{lccc}
\toprule
Sampling configuration & \multicolumn{2}{c}{Dispersion} & Distinct values \\
\cmidrule(lr){2-3}
 & Chronos-Small & Chronos-Mini & out of 1000 draws \\
\midrule
\texttt{top\_k} $=50$ \emph{(default)} & \textbf{0.121} & \textbf{0.115} & \textbf{50} \\
\texttt{top\_k} $=200$ & 0.331 & 0.328 & 186--200 \\
\texttt{top\_k} $=1000$ & 0.958 & 0.945 & 275--626 \\
\texttt{top\_k} $=4094$ (full vocabulary) & 1.087 & 1.059 & 276--757 \\
\midrule
temperature $0.5$, at \texttt{top\_k} $=50$ & 0.120 & 0.114 & 50 \\
temperature $2.0$, at \texttt{top\_k} $=50$ & 0.121 & 0.115 & 50 \\
\texttt{top\_p} $=0.9$, at \texttt{top\_k} $=50$ & 0.115 & 0.109 & 43--45 \\
\texttt{top\_p} $=0.99$, at \texttt{top\_k} $=50$ & 0.121 & 0.115 & 50 \\
\bottomrule
\end{tabular}

\end{table}

All three readings hold (Table~\ref{tab:dose_response}). Dispersion recovers
monotonically --- 0.121, 0.331, 0.958, 1.087 for Chronos-Small --- and reaches
its target only when the vocabulary is not truncated at all. One thousand draws
under the default contain exactly 50 distinct values in all 1600 model--asset--date
cells, without exception; this is a signature of the mechanism rather than
evidence about the model, and it is the reading that makes the diagnosis
arithmetic rather than inferential.

The third row of the lower block is the reason the defect is durable. A fourfold
change in temperature moves dispersion by 0.001, and nucleus sampling moves it by
0.006. These are the parameters that are documented, discussed in the literature
and exposed in every tutorial; an analyst who conducts a careful sensitivity
analysis over them obtains a flat response and correctly reports the model as
robust. The one parameter that matters is the one nobody varies.

\subsection{Reading the distribution instead of sampling it}
\label{sec:mechanism:analytic}

At a one-step horizon the sampler is unnecessary. The predictive law is
categorical and exactly available, so its quantiles can be read in closed form:

\begin{enumerate}
	\item apply the softmax to the logits, giving probabilities over token
	identifiers;
	\item map each identifier to its bin centre, offset by the special-token
	count;
	\item multiply by the tokenizer's own scale, the mean absolute value of the
	context;
	\item sort the support and take the cumulative sum of the probabilities;
	\item read $\mathrm{VaR}_\alpha$ as the smallest value whose cumulative
	probability reaches $\alpha$.
\end{enumerate}

This estimator has no $\texttt{num\_samples}$, no $\texttt{top\_k}$, no
$\texttt{top\_p}$, no temperature and no seed. There is no configuration to get
wrong, and it costs one forward pass.

A reconstruction of a model's support is a claim about that model and is checked
against it before anything is concluded from it. Over 40 dates on SP500 we
compare the analytic quantiles against sampling at full vocabulary with 4000
draws: predictive standard deviations agree to within 0.3\% (1.041 against
1.038 in units of realised volatility), and the violation rates at all four
levels agree to four decimal places. The two routes disagree only where Monte
Carlo error predicts they should. The sampled default, on the same dates, gives
dispersion 0.128 and a 1\% violation rate of 0.3942.

Cost is not a consideration either. Across 24 assets and both checkpoints the
analytic route required one forward pass per date --- 121{,}923 dates per
checkpoint --- and completed in roughly ten minutes per checkpoint on a laptop
GPU, against 1000 sampled paths per date for the sampling route. At horizon one,
reading the distribution is both more accurate than sampling it and cheaper.

\subsection{What this does to the forecasts}
\label{sec:mechanism:consequence}

\begin{table}[htbp]
	\centering
	\caption{Raw violation rates by nominal level, 24 assets, no recalibration of
		any kind. The default configuration produces a series whose exceedance
		rate is almost invariant to the level it was asked for.}
	\label{tab:alpha_response}
	% Generated by Quantlets/CO_chronos_sampling/run_alpha_response.py
\begin{tabular}{lcccccccc}
\toprule
& \multicolumn{4}{c}{Violation rate $\hat\pi(\alpha)$} & \multicolumn{4}{c}{Kupiec passes, of 24} \\
\cmidrule(lr){2-5}\cmidrule(lr){6-9}
& $0.01$ & $0.025$ & $0.05$ & $0.10$ & $0.01$ & $0.025$ & $0.05$ & $0.10$ \\
\midrule
Chronos-Small, sampled at the default & 0.3884 & 0.3962 & 0.4044 & 0.4177 & 0 & 0 & 0 & 0 \\
Chronos-Small, read analytically & 0.0175 & 0.0338 & 0.0577 & 0.1036 & 8 & 10 & 14 & 20 \\
Chronos-Mini, sampled at the default & 0.4188 & 0.4261 & 0.4340 & 0.4477 & 0 & 0 & 0 & 0 \\
Chronos-Mini, read analytically & 0.0178 & 0.0333 & 0.0558 & 0.0989 & 8 & 12 & 17 & 22 \\
\bottomrule
\end{tabular}

\end{table}

Table~\ref{tab:alpha_response} reports the raw forecasts on the full panel. Two
things in it deserve separate statements.

\paragraph{The default produces a series that is not a quantile forecast} Read
the first and third rows across, not down. Between the 1\% and the 10\% level the
nominal exceedance probability changes tenfold; the realised rate changes by a
factor of 1.08. A distribution truncated to its fifty most probable bins has so
little spread that its $\alpha$-quantiles are nearly coincident, and the series
built from them responds to the level it was asked for almost not at all. That
property is diagnostic on its own: it needs no benchmark, no test window and no
backtest, and it is the $\alpha$-response check of
Section~\ref{sec:gate}.

\paragraph{The residual is a tail phenomenon, and it shrinks} Read through its
own predictive distribution, Chronos is not well calibrated at the 1\% level:
0.0175 against a nominal 0.0100 is 1.75 times nominal, and unconditional coverage
rejects on 16 of 24 assets. We state this rather than leave it to be found. But
the residual is not a level shift. The ratio $\hat\pi/\alpha$ falls monotonically
--- 1.75, 1.35, 1.15, 1.04 for Chronos-Small and 1.78, 1.33, 1.12, 0.99 for
Chronos-Mini --- so that at $\alpha = 0.10$ the analytic forecasts are
essentially calibrated, passing Kupiec on 20 and 22 of 24 assets. What survives
the correction is a specifically \emph{tail} deficiency, which is a claim about
the model and its quantisation grid rather than about the interface through which
it is read.



%══════════════════════════════════════════════════════════════════════════════
%─── PART 8: Empirical Results ──────────────────────────────────────────────
%══════════════════════════════════════════════════════════════════════════════
\section{Empirical Results}
\label{sec:results}

\subsection{Audit trail}
\label{sec:audit_trail}

This paper corrects a working paper of our own, and the corrections are
substantial enough that a reader is entitled to see them in one place rather
than assembled from remarks scattered through the text.
Table~\ref{tab:audit_trail} is that place. The rest of the paper states the
current position positively and refers back here; the one exception is
Section~\ref{sec:blindspot_retro}, where what the superseded tables showed is
itself the evidence and has to be read in context.

Each row names the diagnostic that established the defect, because a claim that
a number was wrong is worth no more than the test that shows it.

\begin{table}[htbp]
	\centering
	\caption{Every correction to the working paper, with the diagnostic that
		established it. Rates are raw violation rates at $\alpha = 0.01$,
		averaged over \nMainAssets{} assets.}
	\label{tab:audit_trail}
	\footnotesize
	\setlength{\tabcolsep}{4pt}
	\begin{tabular}{@{}p{3.1cm}p{4.6cm}p{2.4cm}p{2.9cm}@{}}
		\toprule
		\textbf{Defect or claim} & \textbf{Diagnostic} &
		\textbf{As reported} & \textbf{Corrected} \\
		\midrule
		\multicolumn{4}{@{}l}{\emph{Defects in the forecast pipeline}}\\[2pt]
		Sign inversion, TimesFM~2.5 & stored series reproduced by
		$-F^{-1}(\alpha)$ to ten decimals from parameters in the same file
		& $\hat\pi = 0.990$ & $\hat\pi = \nRawPiTimesFM$ \\
		Sign inversion, Moirai~2.0 & as above; 8 of 10 forecast directories
		store the correct sign & $\hat\pi = 0.988$ & $\hat\pi = \nRawPiMoiraiTwo$ \\
		Sampler default \texttt{top\_k} $=50$, both Chronos checkpoints
		& dose--response at fixed context and seed; 1000 draws contain exactly
		50 distinct values & $\hat\pi = \nPiSmallDefaultOne$, $\nPiMiniDefaultOne$
		& $\hat\pi = \nPiSmallAnalyticOne$, $\nPiMiniAnalyticOne$ \\
		Raw $t_5$ quantile with $\nu$ hard-coded, GJR-GARCH & quantile map read
		from the generator; thresholds inflated by 45\%
		& $\hat\pi = 0.004$ & $\hat\pi = \nRawPiGJR$ \\
		\texttt{CACT}/\texttt{FCHI} alias and an undated forecast vintage
		& manifest and ticker audit & entered six scripts
		& quarantined \\[4pt]
		\multicolumn{4}{@{}l}{\emph{Defects in the analysis code}}\\[2pt]
		Sign inversion in both tail-closure scripts & reconstruction compared
		against the promoted series & ``$R>1$ under every closure''
		& $R \in [0.005, 1.70]$ \\
		Sign inversion in the $Z_2$ routine & second implementation written from
		the published definition & 24/24 assets pass at $\hat\pi = 0.39$
		& 0/24; $Z_2 = -144$, $-389$ \\
		Inner-seven-deciles check & fits a Student-$t$ to points generated by a
		Student-$t$ & ``most demanding check'' & cannot fail \\[4pt]
		\multicolumn{4}{@{}l}{\emph{Claims withdrawn or restated}}\\[2pt]
		Predictive interface drives tail failure & corrected panel ordering; the
		two Moirai releases differ on three axes & 1.5\% vs 98.8\% within family
		& \nRawPiMoiraiOne{} vs \nRawPiMoiraiTwo \\
		$R$ as an audit statistic & $\mathrm{Spearman}(\bar R,
		\hat\pi_{\mathrm{raw}}) = \nSpearmanRPi$ & ``dominates both backtests''
		& a re-encoding of $\hat\pi$ \\
		Gating rule cost & re-evaluated on the calibration window
		& half the damage, no upgrades lost
		& \nGateRollAvoidedPct\%, \nGateRollUpgradesLost{} upgrades lost \\
		Share of miscalibration removable by a shift & Gneiting--Resin
		decomposition on the corrected panel & 35\% & \nuMCBSharePct\% \\[4pt]
		\multicolumn{4}{@{}l}{\emph{Repairs to the theory}}\\[2pt]
		Lemma~S.9.1 & the forward $\sigma$-field inclusion fails
		for an infinite context & $\beta_S(k)\le\beta_r(k-1)$
		& $\beta_S(k)\le\beta_r(k-m)$ \\
		Proposition~S.9.4 & the regret bound does not imply the
		coverage statement & regret argument & exact telescoping \\
		Proposition~S.9.6 & drift measure not estimable from
		one path & stated as a bound & labelled heuristic \\
		\bottomrule
	\end{tabular}
\end{table}



The primary evaluation unit is the model--asset pair, because regulatory risk systems are backtested at the desk, asset, or portfolio level. Panel-pooled tests provide complementary evidence of average calibration across the asset universe, not substitutes for per-asset backtesting.

\subsection{\texorpdfstring{Raw Diagnostics and the Ordering by $\bar R$}{Raw Diagnostics and the Ordering by R}}
\label{sec:results_raw}

Table~\ref{tab:master} reports raw and conformally corrected diagnostics across
16 forecasters and 24 assets, ordered by $\bar R$. Read it first for what it does
\emph{not} show. Once the two series sampled at the checkpoint default are set
aside, raw 1\% violation rates run from \nMainRawPiMin{} for CAViaR-AS to \nMainRawPiMax{} for GAS-$t$
--- a factor of three across dynamic-quantile models, score-driven models,
GARCH-family benchmarks, historical simulation and six zero-shot foundation-model
series. Including the two truncated series the range is a factor of 38, and that
factor is a property of a sampling parameter rather than of any forecaster.

The panel is uniformly under-covered at the 1\% level and nothing in it is well
calibrated there. No forecaster passes unconditional coverage on more than 15 of
24 assets, the best being a per-asset CAViaR fitted the ordinary way
\citep{engle2004caviar}; the six foundation-model series pass on between 0 and 13.
Across the whole table \nMainKupiecRawPasses{} of \nMainPairs{} pairs pass raw, and \nMainKupiecCorPasses{} pass after correction.
A reader looking for a foundation model that is competitive at the extreme tail
will not find one here, and a reader looking for the catastrophe reported in
earlier work on this panel will not find that either.

\subsubsection{What the ratio $R$ measures, and what it does not}
\label{sec:R_meaning}

The ratio
\begin{equation}
	\label{eq:replacement_ratio}
	R
	=
	\frac{|\qVstat|}{|\VaR_{\mathrm{raw}}|}
\end{equation}
measures how large the conformal correction is relative to the threshold it
adjusts. Values near zero mean the corrected forecast is the raw forecast; values
above one mean the correction exceeds the raw VaR in magnitude, so the reported
threshold is set by the calibration sample rather than by the model. In this
panel only the two truncated Chronos series exceed one, at $\bar R = 17.3$ and
$23.5$; everything else lies between $0.001$ and $0.36$.

An earlier version of this paper presented $R$ as an audit statistic carrying
information that the standard backtests do not. That claim does not survive its
own test and is withdrawn. Among the well-specified forecasters of this panel,
$R$ is a rank-preserving transform of the raw violation rate: the Spearman
correlation between $\bar R$ and $\hat\pi_{\mathrm{raw}}$ is $\nSpearmanRPi$ across
the \nSpearmanRPiN{} well-specified forecasters of Table~\ref{tab:master}, and a detection rule built on the violation rate
alone reproduces every sensitivity and specificity figure we had attributed to
$R$. This is the sharper form of a referee's objection, and it is correct: $R$ is
a scale-interpretable restatement of a quantity already printed in the adjacent
column, useful for saying \emph{how much} of a forecast survives correction, and
not an independent instrument.

The relation to the decomposition of \citet{gneiting2023regression} is the same
statement in their language. $\qVstat$ estimates the unconditional
miscalibration component uMCB, the part of the miscalibration that a constant
shift can remove, and that is an identity between estimators rather than an
empirical finding. What is empirical is its size: on well-specified series a
constant shift addresses a median of \nuMCBSharePct\% of total miscalibration, so the
majority of what is wrong with these forecasts is conditional and a scalar
correction cannot reach it.

\subsubsection{The quantile-grid interface, and what is left of the claim}
\label{sec:interface_failure}

The working paper reported raw 1\% violation rates of 98.8\% for Moirai~2.0 and
99.0\% for TimesFM~2.5 and read them as evidence that a coarse quantile grid
cannot deliver a backtestable 1\% VaR. Both figures were a sign inversion in our
own generation code, and the corrected series run at $\hat\pi =
\nRawPiMoiraiTwo$ and $\nRawPiTimesFM$ (Table~\ref{tab:audit_trail}).

No interface claim is supportable on the corrected panel, in either direction.
The best raw foundation-model series at the 1\% level is a quantile-grid model,
TimesFM~2.5 at $\nRawPiTimesFM$; the sample-based Moirai~1.1 is at
$\nRawPiMoiraiOne$, the analytic Chronos series at $\nRawPiChronosSmallA$ and
$\nRawPiChronosMiniA$, the quantile-grid Moirai~2.0 at $\nRawPiMoiraiTwo$, and
the sample-based Lag-Llama at $\nRawPiLagLlama$. The two interfaces do not
separate.

What survives is narrower and is a statement about specification rather than
about models. A nine-decile grid whose lowest point is the 10\% quantile does
not contain a 1\% quantile, so a tail-completion step must supply it and becomes
part of the forecaster. That step is consequential: across three assets the same
grids give $\bar R$ between 0.005 and 1.70 depending on whether the closure is
Student-$t$, Gaussian or linear, and on BTC the choice moves both models from
the Basel yellow zone to the red one (Supplement~S.10). Reporting
``TimesFM~2.5 at the 1\% level'' without naming the closure describes an object
that does not exist.\footnote{Every series in this table has been through the
scripted validation gate of Section~\ref{sec:gate} and
Supplement~S.11. Four do not pass it: the two Chronos series
sampled at the checkpoint default, which fail five checks each, and their
analytic counterparts, which fail the coverage band on one asset apiece. The
first two are retained deliberately, as the measured consequence of the default
rather than as forecasters; a blocked series is reported with its diagnosis, not
silently dropped.}

\subsection{Static Recalibration}
\label{sec:recalib}

The corrected columns of Table~\ref{tab:master} show that marginal coverage is restored across the panel: \nMainGreen{} of \nMainPairs{} model--asset pairs reach the Basel Green Zone (\nMainGreenPct\%). What the ordering by $\bar R$ adds is the \emph{amount} of correction required to get there. For thirteen of the sixteen forecasters it is between $0.001$ and $0.18$ --- the corrected forecast is the raw forecast to two or three decimal places --- while the two series sampled at the checkpoint default sit at $17.3$ and $23.5$, where the corrected forecast is an empirical fallback retaining nothing of the base model's tail estimate. Green-zone compliance after correction is therefore uninformative on its own about whether the corrected forecast carries any of the base model's information: the truncated series reach it as comfortably as the rest.

We deliberately do not partition the panel. The ordering by $\bar R$ has structure at both ends and none in the middle: the largest consecutive ratio, a factor of 84, separates the dynamic-quantile forecasters from the rest near the \emph{bottom} of the range, where their raw coverage is already nominal, and the next, a factor of 48, separates Lag-Llama from the two truncated Chronos series at the top. Neither is a boundary between kinds of forecaster, and no single cutoff is distinguished by the data. Any threshold would inherit that arbitrariness, and the continuous statistic conveys the same information without it.

Conditional coverage is a different matter, and the reason is partly that the test is often unavailable. The Christoffersen independence statistic is undefined whenever no two violations fall on consecutive days, which after correction at $\alpha = 0.01$ is the case for \nMainCCUndefined{} of the \nMainPairs{} pairs of the main panel. On the \nMainCCDefined{} pairs where it is defined, \nMainCCPass{} pass --- \nMainCCPassPct\%. We report the informative subset because counting an undefined test as a pass, as is common, would turn the same evidence into a \nMainCCAsPassPct\% pass rate constructed largely from tables that do not exist. The degeneracy rate is itself a result, and Section~\ref{sec:blindspot} treats it as one: at the 1\% level, on samples of the length regulators use, the standard independence test has no evidence to work with on more than half the panel. Green Zone recovery should therefore be read as restored unconditional coverage and nothing more. The structural reasons why a scalar shift cannot reshape conditional dynamics are in Section~\ref{sec:failure}.

\begin{table}[htbp]
	\centering
	\caption{Forecast audit across 16 forecasters and 24 assets ($\alpha = 0.01$), ordered by $\bar R$. QS: Quantile Score ($\times 10^{4}$, lower is better). $\bar R = |\qVstat|/|\VaR_{\mathrm{raw}}|$, averaged over assets: the size of the conformal correction relative to the threshold it adjusts. Kind records the estimator class. The two Chronos rows without the analytic suffix are the checkpoint-default configuration of Section~\ref{sec:mechanism} and are retained in the table as the measured consequence of that default, not as forecasters. No threshold is applied and no classification is implied. $\bar R$ is a ratio of absolute values and is therefore non-negative by construction; the sign that carries the diagnosis belongs to $\qVstat$, and the count of assets on which it is negative is given in the last column.}
	\label{tab:master}
	\tiny
	\setlength{\tabcolsep}{3pt}
	\setlength{\tabcolsep}{1pt}
\begin{tabular}{@{}ll rr rr r rr rr rr@{}}
\toprule
& & \multicolumn{2}{c}{$\hat\pi$} & \multicolumn{2}{c}{Kupiec pass}
& CC pass & \multicolumn{2}{c}{QS} & & & & \\
\cmidrule(lr){3-4}\cmidrule(lr){5-6}\cmidrule(lr){8-9}
Forecaster & Interface & Raw & Corr. & Raw & Corr. & Corr.
& Raw & Corr. & Width & Green & $\bar{R}$ & $\qVstat<0$ \\
\midrule
CAViaR-AS & dyn. quant. & .011 & .011 & 15/24 & 15/24 & 9/11 & 4.8 & 4.8 & .036 & 20/24 & 0.001012 & 2 \\
CAViaR-SAV & dyn. quant. & .011 & .011 & 14/24 & 14/24 & 6/11 & 4.8 & 4.8 & .036 & 18/24 & 0.001084 & 2 \\
TimesFM 2.5 & TSFM, grid & .014 & .010 & 13/24 & 20/24 & 8/16 & 5.0 & 5.0 & .037 & 22/24 & 0.09079 & 3 \\
GJR-GARCH-$t$ & parametric & .014 & .010 & 13/24 & 20/24 & 4/7 & 4.7 & 4.7 & .037 & 23/24 & 0.1011 & 1 \\
Moirai 1.1 & TSFM, sample & .015 & .011 & 13/24 & 15/24 & 6/6 & 5.4 & 5.3 & .039 & 21/24 & 0.1056 & 1 \\
Hist.\ Sim. & nonparam. & .016 & .011 & 11/24 & 21/24 & 6/16 & 5.7 & 5.6 & .039 & 22/24 & 0.1088 & 0 \\
Chronos-Sm. (anal.) & TSFM, anal. & .018 & .011 & 8/24 & 19/24 & 8/14 & 5.2 & 5.1 & .035 & 22/24 & 0.145 & 0 \\
Moirai 2.0 & TSFM, grid & .018 & .011 & 8/24 & 17/24 & 7/15 & 5.2 & 5.1 & .036 & 21/24 & 0.1552 & 2 \\
Chronos-Mn. (anal.) & TSFM, anal. & .018 & .011 & 8/24 & 20/24 & 7/15 & 5.2 & 5.0 & .036 & 22/24 & 0.1608 & 0 \\
GARCH-N & parametric & .019 & .010 & 8/24 & 18/24 & 6/12 & 4.9 & 4.7 & .036 & 21/24 & 0.1713 & 0 \\
GJR-GARCH & parametric & .020 & .011 & 7/24 & 17/24 & 5/6 & 4.8 & 4.7 & .036 & 22/24 & 0.1786 & 0 \\
EWMA & parametric & .021 & .011 & 5/24 & 16/24 & 7/10 & 5.1 & 4.9 & .037 & 20/24 & 0.1843 & 0 \\
GAS-$t$ & score-dr. & .033 & .012 & 0/24 & 14/24 & 7/13 & 5.6 & 4.9 & .036 & 19/24 & 0.3209 & 0 \\
Lag-Llama & TSFM, sample & .029 & .009 & 0/24 & 22/24 & 2/2 & 5.9 & 5.2 & .041 & 24/24 & 0.3573 & 0 \\
Chronos-Small (dflt) & TSFM, sample & .388 & .011 & 0/24 & 14/24 & 4/13 & 37.9 & 5.8 & .040 & 19/24 & 17.26 & 0 \\
Chronos-Mini (dflt) & TSFM, sample & .419 & .010 & 0/24 & 11/24 & 5/13 & 41.1 & 5.7 & .041 & 19/24 & 23.54 & 0 \\
\bottomrule
\end{tabular}

	\begin{minipage}{\linewidth}\scriptsize
		\smallskip
		Corr.: after conformal correction. Width: mean corrected VaR threshold. Green: assets reaching the Basel Green Zone after correction (scaled to 250~days). CC pass is reported as passes/defined: the Christoffersen independence test is undefined when no two violations fall on consecutive days, and an undefined test is not counted as a pass. Summary statistics for the table are in the generated note beneath it, so that they cannot drift from the table itself.
	\end{minipage}
\end{table}

Corrected violation rates cluster tightly around $\hat\pi = 0.010$ for every forecaster in the table, which is what the guarantee promises and is exactly why the corrected column cannot rank them. The Quantile Score can: it runs from $4.67 \times 10^{-4}$ for GJR-GARCH-$t$ to $5.83 \times 10^{-4}$ for the truncated Chronos-Small, and the ordering it produces is close to the ordering by $\bar R$, because both are dominated by how far the raw forecast had to move.

Conformal corrections widen VaR for under-conservative forecasters and contract it for over-conservative ones; for the panel, the dominant effect is widening, raising implied capital requirements (Table~\ref{tab:master}, Width/GJR column). Calibration--efficiency trade-offs and the resulting capital arithmetic are discussed in Sections~\ref{sec:failure} and~\ref{sec:economic}.

Figure~\ref{fig:traffic_light} shows asset-level Basel classifications before and after correction; the raw-versus-corrected violation-rate plot and the calibration--efficiency frontier are in Supplement~S.4.

\begin{figure}[htbp]
	\centering
	\includegraphics[width=\textwidth]{fig_traffic_light}
	\caption{Basel Traffic Light heatmap (13 models $\times$ 24 assets, $\alpha = 0.01$). Top: raw forecasts; Bottom: after conformal correction. Green: $\leq 4$ violations per 250 days; Yellow: 5--9; Red: $\geq 10$. Raw: 106/312 Green (34\%), 154 Yellow, 52 Red. Corrected: 278/312 Green (89\%), 34 Yellow, no Red.}
	\label{fig:traffic_light}
\end{figure}


\subsection{When the Correction Helps, and What It Costs}
\label{sec:when_to_apply}

Aggregate improvements can conceal deteriorations. We therefore report the full
distribution of $\Delta\mathrm{QS} = \mathrm{QS}_{\mathrm{raw}} -
\mathrm{QS}_{\mathrm{corrected}}$ over model--asset pairs rather than its mean.

At $\alpha = 0.01$ the correction makes 66 of 312 pairs worse under the
single-split estimator and 174 of 312 worse under the rolling estimator. The
deteriorations are not scattered: they concentrate on pairs whose raw forecast
was already well calibrated, where the correction has nothing to remove and adds
the estimation noise of $\qVstat$. Of the 94 pairs whose raw forecast passes
Kupiec on the test set, the single-split correction degrades 58 (mean
$\Delta\mathrm{QS} = -1.32\%$ of $\mathrm{QS}_{\mathrm{raw}}$, Wilcoxon
$p = 9.3 \times 10^{-5}$) and the rolling variant degrades 89, with mean
$-7.35\%$ and $p \approx 10^{-16}$. The rolling estimator re-estimates $\qVstat$
on a 250-day window and is correspondingly noisier, which is why its
deterioration rate is roughly double at every level. The pattern holds at
$\alpha \in \{0.025, 0.05, 0.10\}$, where it applies to a larger share of the
panel because more raw forecasts are already acceptable at those levels.

This is a sharper answer than the earlier version of this paper could give. With
two series running at $\hat\pi \approx 0.99$ in the panel, ``deterioration''
was not a meaningful notion for them and their magnitudes dominated every pooled
statistic; removing them makes the effect on the remaining pairs visible rather
than changing it.

Whether a deterioration matters depends on whether it buys anything. Of the 174
pairs the rolling estimator degrades, 72 move from Yellow to Green --- a trade a
risk manager can accept --- but 100 gain no zone change at all, 99 of them
because they were already Green, and 2 move from Green to Yellow, worsening on
both axes. On 59\% of the degradations the correction is therefore not a
trade-off but the wrong intervention.

This suggests a deployment rule that requires no classification: apply the
correction when the raw forecast fails a backtest (Basel zone worse than Green,
or Kupiec rejected), and not otherwise. \emph{Which} backtest matters, and the
distinction is not cosmetic.

Gated on the evaluation window, the rule looks excellent: at $\alpha = 0.01$ it
applies to \nGateRollOracleApplied{} of \nSeqPairs{} pairs, skips
\nGateRollOracleSkipped, avoids \nGateRollOracleAvoided{} of the
\nDegradedRoll{} rolling deteriorations --- \nGateRollOracleAvoidedPct\% of the
damage --- and retains \emph{all} \nGateRollOracleUpgrades{} Basel zone
upgrades. But that rule reads the window it is then scored on. It is an oracle,
and we report it only as an upper bound on what gating can buy.

The deployable rule gates on the calibration window, which is the information
available on the day the decision is taken. It is materially weaker. At
$\alpha = 0.01$ it applies to \nGateRollApplied{} pairs and skips only
\nGateRollSkipped; it avoids \nGateRollAvoided{} deteriorations,
\nGateRollAvoidedPct\% of the damage rather than half; and it costs
\nGateRollUpgradesLost{} of the \nGateRollUpgrades{} zone upgrades, so the claim
that gating is free on the regulatory axis does not survive contact with the
constraint. Under the single split the deployable rule avoids
\nGateStaticAvoided{} of the \nDegradedStatic{} deteriorations and costs
\nGateStaticUpgradesLost{} of \nGateStaticUpgrades{} upgrades.

The honest summary is therefore weaker than the one an in-sample evaluation
would support, and still useful: a calibration-window backtest carries enough
signal to avoid about a quarter of the damage at a cost of about a tenth of the
zone upgrades, and the gap between that and the oracle is a measure of how
poorly the calibration window predicts the test window at $\alpha = 0.01$ ---
which is the tail-sparsity constraint of Remark~\ref{rem:tail_sparsity} showing
up in a decision problem rather than in an estimator. The rule does not bind at
$\alpha \geq 0.025$ in the way an earlier draft claimed. The raw rejection rate
falls steadily with the level --- \nSeqKupiecRejRawOne\%,
\nSeqKupiecRejRawTwoFive\%, \nSeqKupiecRejRawFive\% and
\nSeqKupiecRejRawTen\% at the four levels --- so at $\alpha = 0.10$ a clear
majority of raw forecasts \emph{pass} and the rule skips most of the panel. It
binds at every level, on a growing share of pairs as $\alpha$ rises; what is
specific to $\alpha = 0.01$ is not how often it fires but the size of the damage
it avoids when it does. The decision problem is therefore
specific to the extreme tail --- which is where the identification constraint of
Remark~\ref{rem:tail_sparsity} operates.

\subsection{Rolling recalibration}
\label{sec:rolling_results}

The rolling variant re-estimates $\qVstat$ on a trailing 250-day window, so
$R_t = |\qVroll|/|\widehat{\VaR}^{\mathrm{raw}}_t|$ becomes a monitored series
rather than a single number: sustained elevation says the correction is setting
the reported threshold, transient spikes are stress episodes. No threshold and no
persistence rule are imposed.

It improves coverage substantially and conditional adequacy not at all. Basel
green rises from 89\% (278 of \nSeqPairs) to 99\% (309), with twelve of the
\nSeqForecasters{} forecasters reaching all \nMainAssets{} assets; the exception
is Historical Simulation, whose raw threshold is an order statistic of the same
trailing window the correction uses, so the two estimate from overlapping
information. The gains concentrate where volatility moves: commodities go from
67.3\% green to 100\%, bonds from 89.7\% to 100\%, while crypto and FX are at
100\% either way. Christoffersen pass rates are unchanged within noise. Full
comparison in Table~S.22, the S\&P~500 threshold path in Figure~S.4.

\subsection{Static versus Rolling}
\label{sec:static_vs_rolling}

The rolling variant improves extreme-quantile calibration substantially: Basel Green Zone rates rise from 89\% (278/312) to 99\% (309/312) (Table~S.22). Twelve of the thirteen forecasters reach 24/24 assets Green under the rolling estimator. The exception is Historical Simulation, which falls from 22 to 21 --- the one series whose raw threshold is an order statistic of the same trailing window the rolling correction uses, so the two are estimating from overlapping information.

Conditional coverage does not improve in parallel: GARCH-N achieves the highest rolling Christoffersen pass rate (10/24), while high-$\bar R$ TSFMs remain low (0--5/24), reflecting persistent violation clustering after conformal correction. This is a structural limitation of scalar recalibration, not a tuning issue (Section~\ref{sec:failure}).



Historical Simulation is the only model that deteriorates under rolling ($\hat\pi = 0.013$, 21/24 Green versus 22/24 under single-split), because HS is itself a rolling empirical quantile and an additional rolling conformal layer applies a second window to the same data, reducing responsiveness.

\subsection{Multiple levels, and forecast comparison}
\label{sec:multiquantile}

Evaluated at four levels, the panel separates the two Chronos configurations more
sharply than the 1\% table does, and it does so in the rejection counts rather
than in the violation rates. Corrected rates sit near nominal at every level for
every series --- that is what the guarantee delivers, and it is why the corrected
column cannot rank anything. The per-asset Kupiec counts can. Read analytically,
Chronos-Small is rejected on 5, 9, 5 and 4 of \nMainAssets{} assets across
$\alpha \in \{0.01, 0.025, 0.05, 0.10\}$; sampled at the checkpoint default, the
same weights give 10, 14, 16 and 15. The default configuration is worse at the
off-calibration levels, where a shift calibrated at 1\% does not help it, and the
gap widens as $\alpha$ rises. Recalibration at one level cannot repair a
predictive distribution that is too narrow at all of them (Table~S.12).

At the panel level, EWMA is the only forecaster rejected by the cluster-robust
test ($p = 0.035$). Three others are rejected by the pooled unconditional test
and not by the cluster-robust one --- Moirai~2.0, Chronos-Small read
analytically, and GJR-GARCH, at pooled rates of $0.0114$, $0.0113$ and $0.0111$
against a nominal $0.01$. The pooled test treats 36{,}000 observations as
independent and rejects deviations that are economically negligible; allowing
dependence within an asset removes the rejections (Table~S.13).

Diebold--Mariano comparisons on the corrected forecasts run in one direction:
every significant benchmark--TSFM difference favours the benchmark, 18 of 30 at
the 5\% level, with GJR-GARCH-$t$ the strongest single forecaster on score
(Table~S.23). Historical Simulation is the exception in the other direction,
losing to every TSFM series at $p < 0.001$ --- order-statistic variance on a
250-day window, not a property of the foundation models. The comparison that
belongs here rather than in a robustness appendix is the configuration one: on
identical weights, contexts and dates, reading the Chronos predictive
distribution analytically beats sampling it at the checkpoint default, mean
corrected quantile score $5.03$ against $5.81 \times 10^{-4}$ for Chronos-Small
and $5.00$ against $5.69$ for Chronos-Mini, both at $p < 0.001$. These are the
\emph{corrected} series, so the difference survives a recalibration that has
already restored coverage for both.

\subsection{The Moirai within-family comparison, withdrawn}
\label{sec:interface_vs_architecture}

The working paper presented Moirai~1.1 against Moirai~2.0 as a within-family
control isolating the predictive interface. It is withdrawn for two independent
reasons. The measurement was wrong (Table~\ref{tab:audit_trail}): recomputed, the
two releases run at $\nRawPiMoiraiOne$ and $\nRawPiMoiraiTwo$, a gap of 0.24
percentage points. And the design was not a control: Moirai~1.1 is a
masked-encoder model with multi-patch inputs and a mixture output, while
Moirai~2.0 replaces all three with a decoder-only architecture, a single patch
and a quantile loss, trained on a different corpus
\citep{woo2024moirai,liu2025moirai}. Output format, architecture and pretraining
data change together, so even a large gap could not have been attributed to the
interface.

The second reason is the more useful one to carry forward. A within-family label
is not evidence of a controlled comparison, and what differs between two releases
is a primary-source question to settle before the comparison is drawn rather
than after.

\FloatBarrier
\subsection{Alternative recalibration methods}
\label{sec:baselines}

Ten alternatives to the one-parameter shift were recomputed on the corrected
panel: five single-shot post-hoc corrections, adaptive conformal inference,
gradient-boosted quantile regression, GAMLSS-SST, EVT-POT and filtered historical
simulation (Table~S.24, specifications in Supplement~S.2). Performance is
governed less by flexibility than by whether the recalibration layer respects the
tail-sparsity constraint of Remark~\ref{rem:tail_sparsity}.

Among single-shot procedures the conformal shift has the best quantile score
($5.14 \times 10^{-4}$) and the highest green rate ($89.1\%$); quantile
regression on residuals is close on score at lower coverage; ACI reaches higher
coverage ($97.4\%$) with wider intervals and, at the Basel operating point, a
mean $|\VaR|$ of $0.0750$ against $0.0443$ for the rolling estimator. More
flexible methods do not consistently improve on it: GBM-QR reaches only $50.6\%$
green with a $59\%$ Kupiec rejection rate, and a tuning grid confirms the
prediction that the score-optimal setting buys score at the cost of coverage
(Supplement~S.2.2). GAMLSS-SST is close to the conformal shift on both axes.
Purpose-built VaR models face the same short-window constraint: under the 250-day
protocol, GARCH-N with conformal correction outperforms EVT-POT and filtered
historical simulation, with EVT-POT weakened by GPD estimation failures.

One number in that table is worth isolating. The raw row's mean violation rate is
$0.078$, which is not a statement about forecasters but about the two series
sampled at the checkpoint default; before the corrections of
Section~\ref{sec:interface_failure} it read $0.319$.

\subsection{Simulation and robustness}
\label{sec:simulation}

A Monte Carlo study over five data-generating processes
($T \in \{1{,}000, 5{,}000\}$, 500 replications) confirms the direction of
Theorem~\ref{thm:dependent_conformal_var}: corrected coverage converges toward
nominal as $T$ grows, with Basel green frequencies rising from 76--81\% at
$T = 1{,}000$ to 96--98\% at $T = 5{,}000$ (Table~S.25, design in
Supplement~S.4.1).

Five robustness checks are reported in full in Supplement~S.4. Calibration-fraction
sensitivity is limited over $f_c \in [0.50, 0.80]$; rolling-window length
($w \in \{125, 250, 500\}$) preserves the ordering by $\bar R$; the theoretical
gap $g_n$ has negligible empirical effect, at most $0.0058$ in $\hat\pi$ within
the COVID window and $0.0005$ over the full sample; and block-bootstrap intervals
for $\qVstat$ are narrow relative to the spread across forecasters. The
tail-completion rule is the exception and is not a sensitivity in the usual
sense: it moves $\bar R$ by a factor of 3 to 76 within a single model--asset pair
and can change the Basel zone, which is why Section~\ref{sec:interface_failure}
treats the closure as part of the forecaster rather than as an implementation
detail.

%══════════════════════════════════════════════════════════════════════════════
%─── PART 7b: What the diagnostics could not see ────────────────────────────
%══════════════════════════════════════════════════════════════════════════════
\section{What the Standard Diagnostics Cannot Diagnose}
\label{sec:blindspot}

Four defective series passed every backtest reported in the earlier version of
this paper. This section asks why. The answer is not that the tests were applied
carelessly, and it is not that they are blind to miscalibration --- Kupiec's test
rejects most of this panel, defective series included. It is narrower and more
awkward: a coverage test reports \emph{that} a series miscovers, never
\emph{why}. It cannot separate a model that is genuinely poor in the tail from
one whose predictive distribution was truncated by a sampling default, whose
quantile was stored with the wrong sign, or whose returns were misaligned. All
four present as the same rejection. Add that at $\alpha = 0.01$ these tests have
their least power exactly where the regulatory framework places them, and that
the correction applied downstream of them removes the evidence they would need,
and a defect can pass through a pipeline untouched.

\subsection{The independence test is frequently undefined}
\label{sec:blindspot_cc}

The independence component of \citet{christoffersen1998} conditions on the
transition counts of the exceedance indicator. When no two violations fall on
consecutive days the table is degenerate: $\hat\pi_{11}$ is estimated at the
boundary of the parameter space, the likelihood ratio can be defined by
continuity but its $\chi^2_1$ calibration is not valid there, and the test has
essentially no power against the alternative it is meant to detect. Standard
implementations differ in what they do about this --- ours returns a missing
value, others return a $p$-value computed from a statistic whose asymptotics do
not hold. We call the test undefined in that state, and the substantive point
does not depend on the word: in a degenerate table the test carries no evidence
either way. At $\alpha = 0.01$ this is not a corner case: across our 312
model--asset pairs the test is undefined on 34.6\% of the raw series, and on
53.5\% of the same pairs after recalibration. The share falls monotonically with
the nominal level --- 11.5\%, 1.0\% and 0.0\% raw at $\alpha = 0.025$, $0.05$ and
$0.10$ --- so this is a property of the tail probability, not of our panel.

Two things follow. First, recalibration makes the problem worse rather than
better, because a correction that restores coverage also spaces violations out,
and the tests are usually run on the corrected series. Second, and more
seriously, a degenerate table is an absence of evidence, but the natural
implementation returns ``no rejection'', which reads as a pass. A test with no power on half of its
applications, reported as passing on all of them, is not evidence of
independence. Where the test \emph{is} defined it
rejects on 27.5\% of raw pairs and 48.3\% of recalibrated ones, so the
non-degenerate cases are not reassuring either.

\subsection{Unconditional coverage rejects almost everything}
\label{sec:blindspot_kupiec}

The test of \citet{kupiec1995} is always defined, and at $\alpha = 0.01$ it
rejects 69.9\% of the 312 pairs, against 39.4\% at $\alpha = 0.10$. No forecaster
in the panel passes on more than \nMainBestKupiec{} of \nMainAssets{} assets at the 1\% level. A test that
rejects the majority of defensible models carries little information about any
one of them, and --- decisively for the question of this paper --- it does not
separate a series with a traced defect from one that is merely imperfect.

We measured that separation rather than asserting it. Treating the five series
with traced defects as positives and the eight without as negatives, Kupiec's
majority-fail rule attains sensitivity 5/5 at specificity 2/8, a Fisher exact
$p$ of $0.49$: it flags 11 of 13 forecasters, which is not discrimination. The
independence test attains 0/5. A severe cut on the violation rate --- rejecting
$\hat\pi > 0.10$ or $\hat\pi < 0.005$ --- attains 5/5 at 8/8, $p = 0.00078$. The
finding is about a \emph{threshold}, not about a statistic: what separates defects
from ordinary miscalibration is severity, and Kupiec's threshold is calibrated to
sample size instead. All six Kupiec flags on forecasters with no traced defect
were followed to a verdict, as pre-registered; none is a defect, so the 2/8
figure stands as measured rather than as an artefact of unexamined positives.

\subsection{Recalibration conceals what it corrects}
\label{sec:blindspot_conformal}

Split-conformal recalibration targets marginal coverage regardless of the
forecaster it is applied to --- exactly in finite samples when the scores are
exchangeable \citep{vovk2005algorithmic,lei2018distribution}, and asymptotically
under the dependence assumptions of Theorem~\ref{thm:dependent_conformal_var},
which our contiguous split does not satisfy exactly
(Remark~\ref{rem:gap_vs_protocol}). What matters for this section is what both
versions share: the guarantee is indifferent to the forecaster. In this panel it
delivers:
335 of 384 pairs reach the Basel green zone after correction. Applied to the
Chronos series truncated at $\texttt{top\_k} = 50$ --- a series whose predictive
dispersion is 0.12 of realised volatility and whose exceedance rate is nearly
invariant to the nominal level --- it produces $\hat\pi = 0.0108$ and places 19 of
24 assets in the green zone. Applied to the sign-inverted series it did the same.

Coverage-based backtests downstream of that step therefore lose most of their
power to separate: after correction \nMainGreen{} of \nMainPairs{} pairs are in
the green zone and the violation rates of every forecaster cluster within a
fraction of a percentage point of nominal, defective series included. This
is not a failure of the method: a correction that guarantees coverage cannot
diagnose the forecaster it corrects, precisely because the guarantee holds
whatever the forecaster is. The practical implication is that a pipeline which
validates \emph{after} recalibration has, by construction, no test of the
forecaster at all.

\subsection{The uncomfortable retrospective}
\label{sec:blindspot_retro}

The instruments were not silent. In the submitted version of this paper, Table~1
partitioned the panel on $R > 1$ into ``signal-preserving'' and ``effective
replacement'' regimes. The replacement panel contained exactly four series ---
Chronos-Small at $R = 17.3$, Chronos-Mini at $23.5$, TimesFM~2.5 and Moirai~2.0
at $3.2$ --- and those four are exactly the four defective foundation-model
series. The fifth defect appears in the same submission: in the cross-sectional
table, GJR-GARCH's correlation between $\qVstat$ and annualised volatility is
$-0.786$, the only negative value against $+0.67$ to $+0.97$ for every other
model, and the manuscript explained it as conditional variance dynamics absorbing
volatility into raw VaR.

Every defect family was flagged, in printed tables, before anyone knew a defect
existed --- and every flag was read as a property of the model. This disposes of
the obvious objection that the detection result above is circular, since the
flags predate the discovery. It also states the real lesson more sharply than any
sensitivity figure can: the failure was one of interpretation, not of
instrumentation. A large correction was reported as a finding about a forecaster
when it should have been read as a question about a pipeline.


%══════════════════════════════════════════════════════════════════════════════
%─── PART 7c: Structural validation ─────────────────────────────────────────
%══════════════════════════════════════════════════════════════════════════════
\section{Structural Validation}
\label{sec:gate}

If coverage cannot detect these failures, validation must not be based on
coverage. We therefore check a forecast series against properties it must have
\emph{to be a Value-at-Risk series at all}, before it is calibrated, evaluated or
compared with anything.

Ten such checks are stated in Supplement~S.11 and implemented
in \texttt{scripts/promotion\_gate.py}: a lower quantile whose median is not
negative, quantiles that are not ordered across $\alpha$, an exceedance rate that
does not respond to the nominal level, a predictive dispersion far from realised
volatility, a support of a handful of distinct values.

They are not all the same kind of thing, and we distinguish them rather than
present ten items of equal standing. Two are invariants that follow from what a
quantile forecast is and carry no tunable constant; one is an invariant tested
through a proxy with a threshold; the remaining seven are plausibility bands
whose cutoffs we chose and state. Eight of the ten need nothing but the forecast
series and the returns. The other two --- $\alpha$-responsiveness and coverage
plausibility --- need an evaluation window, so the gate is not entirely upstream
of the data a backtest consumes. What separates them from a backtest is the
question they ask: not whether a coverage rate is statistically distinguishable
from nominal, but whether the series responds to the level it was asked for at
all. No $p$-value is computed, no benchmark is required, and none of the ten
needs access to the hardware that produced the forecasts.

\paragraph{Scoping is declared, not assumed} Two checks are inapplicable to
order-statistic estimators, which may legitimately tie across adjacent $\alpha$
and hold constant for long stretches as a window slides; a discrete predictive
support may tie for the same reason. The gate records the estimator class of each
series in a manifest and declares the check inapplicable rather than returning a
verdict on it. A check that cannot fail informatively should say so.

\begin{table}[htbp]
	\centering
	\caption{The ten checks, stated on the lower quantile $\hat q^{lo}_\alpha$
		(negative), which is what the forecast files store. Full conditions,
		scoping rules and the per-series results are in Supplement~S.11.}
	\label{tab:gate_compact}
	\footnotesize
	\begin{tabular}{@{}llp{5.4cm}@{}}
		\toprule
		\textbf{Check} & \textbf{Kind} & \textbf{Needs} \\
		\midrule
		Sign & invariant & forecast only \\
		Monotonicity across $\alpha$ & invariant & forecast only \\
		Alignment (no look-ahead) & invariant, proxy test & returns \\
		\midrule
		Support cardinality & band & forecast only \\
		Extremes & band & forecast only \\
		Scale against realised $\sigma$ & band & realised $\sigma$ \\
		Predictive dispersion & band & realised $\sigma$ \\
		$\alpha$-responsiveness & band & evaluation window \\
		Coverage plausibility & band & evaluation window \\
		Tail reach & band & sample paths \\
		\bottomrule
	\end{tabular}
\end{table}

\paragraph{Result on this panel} The gate blocks \nGateBlocked{} of the \nGateSeries{} promoted series.
Two are the Chronos series sampled at the checkpoint default, which fail
monotonicity, scale, $\alpha$-response, coverage and dispersion on all 24 assets.
Two are their analytic counterparts, which fail only the coverage-plausibility
band, and only on one asset each --- the residual tail deficiency of
Section~\ref{sec:mechanism:consequence}, correctly recorded as a block requiring
a written diagnosis rather than silently admitted.

\paragraph{What this evidence is worth} These checks were written after the
defects were known, several of them because a specific defect had already got
past whatever checks existed at the time. Their performance on those defects is
in-sample by construction and is not evidence of their general sensitivity, and
we do not present it as such. What can be claimed is narrower and still useful:
each check is a property a practitioner would endorse before seeing any data,
each is arithmetic rather than inferential, and each addresses a failure mode
that the coverage-based tests of Section~\ref{sec:blindspot} provably cannot
address. The truncated series fail five of the ten checks; the analytic series
fail none of the structural ones.



	%══════════════════════════════════════════════════════════════════════════════
	%─── PART 9: Discussion ─────────────────────────────────────────────────────
	%══════════════════════════════════════════════════════════════════════════════
\section{Discussion}
\label{sec:discuss}

\subsection{Methodological Limitations}
\label{sec:failure}

As documented in Section~\ref{sec:recalib}, a scalar shift restores unconditional violation frequency but cannot reshape the conditional dynamics of the violation process; this is a structural limitation, not a tuning issue. The rolling variant captures time variation in miscalibration magnitude but remains a location adjustment, not a fully conditional recalibration rule. More flexible state-dependent recalibration, such as one-sided extensions of Adaptive Conformal Inference \citep{gibbs2021adaptive} or the PID framework of \citet{angelopoulos2024pid}, lies beyond the present framework.

A related limitation concerns identification. The positive cross-sectional association between $\qVstat$ and asset-level volatility or tail frequency (Supplement~S.5.1) is consistent with pre-training domain mismatch, but the zero-shot design does not permit causal attribution. The audit identifies \emph{where} tail signal is missing, not \emph{why}.

\subsection{Operational Failure Modes}
\label{sec:operational_failure}

Operational failure modes can be mitigated. Small calibration sets ($n_{\mathrm{cal}} < 200$) produce unstable corrections; a minimum calibration size or cross-asset aggregation is advisable. Very high score persistence ($\hat\rho > 0.95$) weakens the $\beta$-mixing bound, and GARCH pre-filtering can reduce dependence. Over-conservative base forecasts ($\qVstat < 0$) may lose coverage under non-stationarity; imposing $\qVstat \geq 0$ provides a guard for systemically over-conservative models. Structural breaks can cause temporary under-protection when the calibration window predates the break; shortening $w$ or incorporating stressed subsamples increases adaptivity at the cost of higher variance.

Figure~S.5 illustrates the structural-break case during the COVID-19 shock. The trailing-window design produces a measurable lag: seven of the thirteen forecasters reach their post-shock $\qVroll$ maximum within the acute window, four peak 77 calendar days later, and the two GJR-GARCH variants take 161 days. The lag is a property of the window rather than of the forecaster class --- Historical Simulation and EWMA are among the fastest, GARCH-N among the slowest. The drift diagnostic $\hat{\delta}_w(t)$ flags this regime in real time (Supplement~S.4, Figure~S.1 of the supplement), allowing operational deployment to detect the lag before it propagates.



\subsection{Capital arithmetic, and what it now shows}
\label{sec:economic}

Under the Basel~II/pre-FRTB internal-models framework the traffic-light zone sets
a capital multiplier $k \in [3,4]$. With $W$ the widening of the corrected
threshold relative to an uncorrected benchmark, the capital ratio between a
corrected Green-zone forecast and an uncorrected Yellow-zone one is
$(k_{\mathrm{Green}}/k_{\mathrm{Yellow}})\,W$, so correction pays whenever
$W < k_{\mathrm{Yellow}}/k_{\mathrm{Green}} \approx 1.22$. Under FRTB the binding
measure is ES at 97.5\% and this arithmetic describes the backtesting channel
rather than the capital channel; we report it as illustrative.

On the corrected panel nothing comes close to the break-even widening. Relative
to GJR-GARCH the corrected thresholds run from $W = 0.98$ for the analytic
Chronos series to $W = 1.14$ for Lag-Llama, with TimesFM~2.5 at
$\nWidthRatioTimesFM$ and Moirai~2.0 at $\nWidthRatioMoiraiTwo$. Correction is
therefore capital-cheap across the board --- a duller statement than the one this
subsection used to make, and a true one. The earlier $W \approx 1.6$--$1.75$ for
the two quantile-grid models was the sign defect, whose inverted thresholds the
correction had to replace outright.

\subsection{Practitioner guidance}
\label{sec:practitioner}

\emph{Validate before you calibrate.} A recalibration decision taken on a series
that has not been structurally checked is a decision about an artefact, and the
correction will make it look correct (Section~\ref{sec:gate}).

\emph{Then decide whether to correct at all.} Apply the shift when the raw
forecast fails a backtest \emph{on the calibration window}, and not otherwise. At
$\alpha = 0.01$ that rule skips \nGateRollSkipped{} of \nSeqPairs{} pairs, avoids
\nGateRollAvoided{} of the \nDegradedRoll{} degradations the rolling estimator
would otherwise cause, and costs \nGateRollUpgradesLost{} of
\nGateRollUpgrades{} Basel zone upgrades. It is not free on the regulatory axis,
and a version that appears to be free is one evaluated on the window it used to
decide.

\emph{Choose the forecaster for the constraint, not the label.} With adequate
per-asset history, GJR-GARCH-$t$ with conformal correction has the best corrected
quantile score in this panel. Where per-asset estimation is infeasible --- scale,
short history, latency --- a correctly configured zero-shot foundation model with
the same correction is competitive with the GARCH family at the 1\% level, which
is a real operational option and not a claim of superiority.

\emph{Monitor the signed shift, not the ratio.} $\qVstat$ carries the diagnosis:
negative means an over-conservative base that the correction tightens, releasing
capital. $\bar R$ cannot carry it, being a ratio of absolute values. On the
corrected panel negative values are rare and belong to no single forecaster ---
three assets for TimesFM~2.5, two for Moirai~2.0, one each for Moirai~1.1 and
GJR-GARCH-$t$. GJR-GARCH itself is now positive on all \nMainAssets{} assets; it
was the systematically over-conservative forecaster of earlier drafts only
because its quantile map inflated every threshold by 45\%.

%══════════════════════════════════════════════════════════════════════════════
\section{Conclusion}
\label{sec:conclude}

A Value-at-Risk forecast is produced by code, and the step in which a predictive
distribution becomes a threshold is where the pipelines we examined failed. In
this panel that step, not the architecture of the model or the format in which it
exposes its distribution, is what separates a forecaster that is roughly usable
at the 1\% level from one that violates on 40\% of days. The clearest instance is
a documented sampling default in a public checkpoint --- correct for the task it
was written for, silently catastrophic for this one --- and the mechanism is
established by controlled experiment rather than inferred: dispersion recovers
monotonically as the truncation is relaxed, one thousand draws under the default
contain exactly fifty distinct values, and the parameters a careful analyst would
vary move nothing.

The finding that matters beyond this panel is that the instruments in standard
use cannot see such failures. At $\alpha = 0.01$ the independence test is
undefined on a third of raw pairs and on more than half of recalibrated ones,
while unconditional coverage rejects seven pairs in ten and so cannot separate a
defect from ordinary miscalibration. Post-hoc recalibration compounds the
problem, because a guarantee that holds regardless of the forecaster cannot
diagnose the forecaster: applied to a truncated series it restored
$\hat\pi = 0.0108$ and 19 of 24 Basel green zones, and every backtest downstream
of it passed. Validation that runs after a correction is not validation of the
model.

We therefore propose checking forecast series structurally, against properties
they must have to be Value-at-Risk series at all, before they are calibrated or
evaluated. Ten such checks are cheap, require no benchmark or test window, and
block 4 of the 13 series promoted here. We state plainly that they were written
after the defects were known and that their performance on those defects is
in-sample; the argument for them is not their measured sensitivity but that the
alternative provably cannot do the job.

On recalibration itself, our claims narrow rather than change. The conformal
shift restores marginal coverage under $\beta$-mixing dependence, and that is
all it does. It is an intervention with an indication: under the rolling
estimator, \nWellCalRollWorse{} of the \nWellCalRollN{} pairs that pass Kupiec on
the test window are degraded by correcting them, and gating the correction on a
calibration-window backtest avoids \nGateRollAvoidedPct\% of that damage at the
cost of \nGateRollUpgradesLost{} of \nGateRollUpgrades{} Basel zone upgrades ---
a quarter of the benefit that the same rule shows when it is allowed to peek at
the evaluation window, which is how such rules are usually reported. It
reaches only part of the problem --- in the decomposition of
\citet{gneiting2023regression} a constant shift addresses a median of \nuMCBSharePct\% of
total miscalibration on well-specified series, the remainder being conditional
--- and it does not restore violation independence. The magnitude $R$ that an
earlier version of this paper offered as an audit statistic is withdrawn as such:
among well-specified forecasters it is a rank-preserving transform of the
violation rate, and we retain it only as a scale-interpretable statement of how
much of a forecast survives correction.

What the corrected panel says about foundation models is deflationary and, we
think, more useful than what we previously reported. Configured correctly, the
zero-shot series here are neither the catastrophe of our earlier draft nor
competitive: they sit alongside GARCH-family benchmarks at the 1\% level, a
per-asset CAViaR fitted the ordinary way outperforms all of them, and the
residual miscalibration of the best of them is concentrated in the tail and
largely gone by $\alpha = 0.10$. Whether fine-tuning closes that tail gap remains
open. What is no longer open, in our view, is whether a tail-risk evaluation can
be trusted without a structural check of the series it evaluates.

\section*{Data and Code Availability}

The full replication pipeline --- data retrieval, foundation-model inference, conformal calibration, backtesting, robustness analysis, Monte Carlo simulation, and figure generation --- together with machine-readable result matrices for all 384 model--asset pairs ($\qVstat$, $\hat\pi$, Basel Traffic Light, Kupiec $p$, Christoffersen $p$, Quantile Score), is available via the Quantlet repository \quantlet{https://github.com/QuantLet/Conformal_Oracle/} at \url{https://github.com/QuantLet/Conformal_Oracle}.

Return data are retrieved from Yahoo Finance. All foundation model checkpoints are open-weight and hosted on HuggingFace; no proprietary API access is required.

A Python package implementing the audit, \texttt{conformal-oracle}, is available on PyPI (\url{https://pypi.org/project/conformal-oracle/}) and installable via \texttt{pip install conformal-oracle}. \revised{Every rolling-mode result reported in this paper was produced by the Quantlet replication pipeline, which has throughout computed the conformal correction as the finite-sample split-conformal quantile---the $\lceil(w+1)(1-\alpha)\rceil$-th order statistic of the trailing-window scores. Version~0.3.1 of the companion \texttt{conformal-oracle} package aligns its packaged rolling implementation with this same convention, which earlier package versions had implemented as the plain empirical $(1-\alpha)$ quantile; consequently no result reported here depended on the affected code path.}


	\section*{Acknowledgements}
	This paper is supported through the project ``IDA Institute of Digital Assets'', CF166/15.11.2022, CN760046/23.05.2023; the project ``AI for Energy Finance (AI4EFin)'', CF162/15.11.2022, CN760048/23.05.2023, financed under the Romania's National Recovery and Resilience Plan, Call no. PNRR-III-C9-2022-I8; and the Marie Skłodowska-Curie Actions under the European Union's Horizon Europe research and innovation program for the Industrial Doctoral Network on Digital Finance, acronym DIGITAL, Project No. 101119635.
	
	We acknowledge the support of the project ``MA'AT -- Autonomous
	Model for Textual Assistance'', SMIS Code 2021+: 330941, funding contract no.\ 390090/11.11.2025, project co-financed by the European Regional Development Fund through the Smart Growth, Digitalisation and Financial Instruments Programme 2021--2027 (POCIDIF).

	\paragraph*{Declaration of generative AI use}
	The authors used Claude (Anthropic), ChatGPT (OpenAI), and Gemini (Google) for language editing and code review during manuscript preparation. All scientific content, analysis, and conclusions are the authors' own.

% Long DOIs and URLs in the reference list overrun the measure otherwise.
\sloppy\emergencystretch=3em
	\bibliographystyle{elsarticle-harv}
	\bibliography{calibrating_the_oracle}
	
	
	%══════════════════════════════════════════════════════
	%─── APPENDICES ─────────────────────────────────────
	%══════════════════════════════════════════════════════
\newpage


\end{document}